% ================================================================
% CHAPTER: CHIRAL KOSZUL DUALITY - FROM QUADRATIC TO NON-QUADRATIC
% ================================================================

\chapter{Chiral Koszul pairs}
\label{chap:koszul-pair-structure}

\index{Koszul pair!recognition criteria|textbf}
\index{Koszul pair!chiral|textbf}
\index{Koszul reflection!the pair (A, A!) as the K-orbit}

\begin{remark}[The five objects in a Koszul pair]
\label{rem:koszul-pair-K-orbit}
\index{Koszul reflection!pair as orbit}
A chiral Koszul pair starts with an augmented chiral algebra~$\cA$
and a conilpotent dual coalgebra~$\cA^i$.  The bar coalgebra
\[
 \Bbarch_X(\cA)=T^c(s^{-1}\bar\cA)
\]
is defined before any dual object is named; Koszulness asserts that
the canonical comparison $\Bbarch_X(\cA)\to\cA^i$ is a
quasi-isomorphism.  The counit
\[
 \Omegach_X\Bbarch_X(\cA)\longrightarrow \cA
\]
is bar--cobar inversion: it recovers the input algebra~$\cA$.  It
does not define~$\cA^!$.

The Koszul-dual algebra belongs to a separate finite-type/Verdier
lane.  Theorem~\textup{\ref{thm:koszul-reflection}} separates the two
functors
\[
 R_X(\cA)=\Omegach_X\Bbarch_X(\cA),\qquad
 K_X(\cA)=\mathbb{D}_{\Ran}\Bbarch_X(\cA).
\]
The first functor is reconstruction; the second is Verdier-dualized
Koszul duality. The Verdier dual of the bar coalgebra is a
factorisation algebra; on the strict finite-type quadratic surface it
reduces to the algebraic dual $\cA^!=(\cA^i)^\vee$.
The symbol $K$ in Theorem~\textup{\ref{thm:koszul-reflection}}
denotes this induced Verdier/Koszul reflection on the localized
Koszul surface.  It is not the raw bar functor, and the statement
$K^2\simeq\mathrm{id}$ is not a chain-level assertion about curved
CDG objects without the localization and completion hypotheses.

The objects are therefore distinct:
\begin{itemize}
\item $\cA$ is the algebra whose OPE data define the bar differential;
\item $\Bbarch_X(\cA)$ is the conilpotent bar coalgebra and carries
 the universal twisting morphism $\tau\colon\Bbarch_X(\cA)\to\cA$;
\item $\cA^i$ is the Koszul-dual coalgebra resolved by the bar
 coalgebra on the Koszul locus;
\item $\cA^!$ is the finite-type/Verdier dual algebra, not the cobar
 of $\Bbarch_X(\cA)$;
\item $Z^{\mathrm{der}}_{\mathrm{ch}}(\cA)=
 R\!\operatorname{Hom}^{\mathrm{ch}}_{\cA^e}(\cA,\cA)$ is the chiral
 derived centre, the Hochschild closed-sector object, not a bar coalgebra and
 not the Koszul-dual algebra.
\end{itemize}
\end{remark}

A Koszul pair $(\cA,\cA^i)$, together with the Verdier-dual algebra
$\cA^!$ when the finite-type hypotheses supply it, is the compatible
package of bar construction, twisting morphism, curvature data, and
deformation filtration from which the MC element $\Theta_\cA$ is
constructed.  The structural theory has three levels of generality
(classical, chiral, and modular); the recognition criteria determine
when the MC element admits the full complement of projections
(Theorems~A--D).

At the classical level, the algebra $\mathrm{Sym}(V)$ is paired first
with its quadratic dual coalgebra; finite-dimensional linear duality
then gives the exterior algebra $\Lambda(V^*)$.  The pair
$(\mathrm{Sym}(V), \Lambda(V^*))$ is detected by the acyclicity of the
twisted tensor product
$\barB(\mathrm{Sym}(V)) \otimes_\tau \Lambda(V^*)$ and the quadratic
presentation of the relation ideal
(Definition~\ref{def:koszul-pair-classical}). This is a theorem about
algebras on a point, with $m_0 = 0$ and $d^2 = 0$ built into the
formalism. On a curve~$X$, the chiral pair
$(\chirCom(V \otimes \omX), \chirLie(V^* \otimes \omX))$ replaces
tensor powers by factorization products on $\FM_n(X)$ and the
classical twisting cochain by a chiral twisting morphism valued in the
convolution algebra on $\Ran(X)$
(Definition~\ref{def:chiral-koszul-pair}, Chapter~\ref{chap:koszul-pairs}).
The genus-zero differential $\dzero$ still satisfies $\dzero^2 = 0$:
no curvature, but an entirely new geometric substrate. When the curve
deforms to genus~$g \geq 1$, the scalar uniform-weight channel
acquires central fiber curvature
$\dfib^2 = \kappa(\cA) \cdot \omega_g$; in multi-weight families this
is the scalar component of a completed CDG object with cross-channel
terms.  The corrected differential $\Dg{g}$ absorbs the curvature into
a curved $\Ainf$-structure with $m_0 \neq 0$ and
$m_1^2(a) = [m_0, a]$. The modular Koszul object
(Definition~\ref{def:modular-koszul-homotopy}) packages all three
levels into a single tower over $\overline{\mathcal{M}}_g$.

The PBW bridge (Theorem~\ref{thm:pbw-koszulness-criterion}) resolves
the central tension: chiral algebras are almost never quadratic, yet
their associated graded under bar-length filtration is, and it is this
graded shadow that carries the classical Koszulness from which the
full chiral pair inherits its structure.


\section{Koszul duality}

\subsection{The classical pair on a point}

Let $V$ be a finite-dimensional vector space. The symmetric algebra
$\mathrm{Sym}(V)$ and the exterior algebra $\Lambda(V^*)$ form a
Koszul pair (Definition~\ref{def:koszul-pair-classical}): there is a
natural twisting cochain
$\tau\colon \barB(\mathrm{Sym}(V)) \to \Lambda(V^*)$ that makes the
twisted tensor product
$\barB(\mathrm{Sym}(V)) \otimes_\tau \Lambda(V^*)$ acyclic. The bar
coalgebra
$\barB(\mathrm{Sym}(V)) = T^c(s^{-1}\overline{\mathrm{Sym}(V)})$
has differential built from the commutative product; its homology is
the quadratic dual coalgebra, and finite-dimensional duality gives the
algebra $\Lambda(V^*)$. Priddy's theorem \cite{Priddy70} gives the
counit
\[
 \Omega(\barB(\mathrm{Sym}(V))) \xrightarrow{\;\sim\;} \mathrm{Sym}(V)
\]
as a quasi-isomorphism. The curvature vanishes: $m_0 = 0$, so the bar
differential satisfies $d^2 = 0$ strictly. All of this is a
shadow of the $\mathrm{Com}^! = \mathrm{Lie}$ operadic duality.

\subsection{The pair on a curve}

Now take $X$ a smooth algebraic curve and replace
$\mathrm{Sym}(V)$ by the chiral symmetric algebra
$\chirCom(V \otimes \omega_X)$, a chiral algebra on $X$ whose fiber
at each point is a copy of $\mathrm{Sym}(V)$, but whose
$n$-point operations are governed by configuration spaces
$\overline{C}_n(X)$ of colliding points rather than by tensor powers.
The quadratic dual coalgebra is represented by the chiral Lie side;
after finite-type dualization this gives the chiral Lie algebra
$\chirLie(V^* \otimes \omega_X)$.  The bar construction
$\barB_X(\cA)$ lives on the Fulton--MacPherson compactification
$\FM_n(X)$. The chiral twisting cochain
(Definition~\ref{def:chiral-koszul-pair}) produces the counit
\[
 \Omega(\barB_X(\cA)) \xrightarrow{\;\sim\;} \cA,
\]
which is again a quasi-isomorphism when $\cA$ is chiral Koszul
(Theorem~A, Theorem~\ref{thm:bar-cobar-isomorphism-main}).  This
map is inversion and has target~$\cA$; the Verdier-dual construction
of~$\cA^!$ is a separate comparison on the bar coalgebra.  The genus-$0$ collision differential $\dzero$ still satisfies
$\dzero^2 = 0$. There is no curvature; the only new ingredient is
the integration kernel mediating chiral operations on $X$.

When $X = \bP^1$, the Beilinson--Drinfeld equivalence identifies chiral algebras on~$\bP^1$ with vertex algebras. The chiral bar construction on~$\bP^1$, however, retains the full configuration-space geometry: the Fulton--MacPherson compactifications~$\FM_n(\bP^1)$, the Arnold relations $\eta_{ij} \wedge \eta_{jk} + \mathrm{cyc.} = 0$, and the logarithmic propagator $\eta_{ij} = d\log(z_i - z_j)$ are all genuinely present even at genus~$0$. The classical bar construction over a point is recovered only upon formal-disk restriction followed by explicit comparison data:
\begin{enumerate}[label=\textup{(\alph*)}]
\item \emph{Formal disk vs.\ point.} Vertex algebras live on the formal disk~$D_x \subset X$, not on a bare point; the passage from~$D_x$ to a point discards the thickening data (completion, growth conditions) that the $\cD_X$-module structure depends on.
\item \emph{Homotopy retract.} The affine line $\bA^1 \hookrightarrow \bP^1$ deformation-retracts onto a point, but this retraction is additional data: a specific chain homotopy and its attendant homotopy transfer (Homological Perturbation Lemma) are required to relate the chiral bar on~$\bA^1$ to the classical bar over the ground field.
\item \emph{Global topology.} On~$\bP^1$ itself, the global topology of configuration spaces (the identification $\overline{C}_{n+1}(\bP^1) \cong \overline{\cM}_{0,n+2}$) introduces boundary strata that have no classical analogue.
\end{enumerate}
The comparison between chiral and classical bar is therefore a quasi-isomorphism mediated by explicit transfer data, not a tautological restriction.

\subsection{The pair in families: curvature from moduli}

Fix a chiral Koszul pair $(\cA,\cA^i)$ on a genus-$0$ curve, with
Verdier-dual algebra~$\cA^!$ only on the finite-type surface, and
deform the curve to genus~$g \geq 1$. The bar construction
$\barB_g(\cA)$ now lives over $\overline{\mathcal{M}}_g$. On the
scalar uniform-weight lane the fiberwise operator has central square
\[
 \dfib^2 = \kappa(\cA)\cdot \omega_g,
\]
where $\omega_g=c_1(\lambda_g)\in H^2(\overline{\mathcal{M}}_g)$ is
the Hodge class and $\kappa(\cA)$ is the modular characteristic
(Theorem~D). If $\kappa(\cA)=0$ the scalar lane is uncurved. If
$\kappa(\cA)\neq0$, $\dfib$ is not an ordinary differential on the
fiber; it is the degree-one part of a CDG object read in the completed
coderived ambient. For multi-weight algebras the same scalar term is
only the uniform-weight projection of the full genus expansion, which
also contains the cross-channel contribution of
Theorem~\ref{thm:multi-weight-genus-expansion}.

The total corrected differential $\Dg{g}$ absorbs the curvature into a higher-order correction, a curved $A_\infty$-structure in which $m_0 \neq 0$ and $m_1^2(a) = [m_0, a]$, and satisfies $\Dg{g}^2 = 0$. The passage from $\dzero$ to $\Dg{g}$ is the passage from classical to quantum: the quantum corrections are controlled by
$H^*(\overline{\mathcal{M}}_g,
Z^{\mathrm{der}}_{\mathrm{ch}}(\cA))$, the cochain-level chiral
closed-sector object.  They are not encoded by $\cA^!$ alone.  Theorem~B
(Theorem~\ref{thm:higher-genus-inversion}) shows that bar-cobar
inversion persists via $E_2$~collapse of the associated spectral
sequence. Theorem~C (Theorem~\ref{thm:quantum-complementarity-main})
identifies the obstruction and deformation spaces as Lagrangian
complements.

\begin{proposition}[Scalar Verdier sums on the finite-type lane;
\ClaimStatusConditional]
\label{prop:kps-scalar-verdier-sums}
\index{complementarity!scalar Verdier sums}
\index{Koszul pair!scalar complementarity constants}
Assume the finite-type Verdier-Koszul comparison exists for the row in
question, and write
\[
 K^\kappa(\cA)=\kappa(\cA)+\kappa(\cA^!).
\]
On the standard witness rows,
\[
\begin{array}{c|c}
\textup{row} & K^\kappa(\cA)\\
\hline
\cH_k,\ V_k(\fg),\ \beta\gamma_\lambda & 0\\
\mathrm{Vir}_c & 13
\end{array}
\]
and on the two $\mathsf M$-extensions and the Vol.~III
$\mathsf B$-row one has
\[
K^\kappa(\mathcal W_3)=\frac{250}{3},\qquad
K^\kappa(\mathrm{BP})=\frac{98}{3},\qquad
K^\kappa(\mathcal H_{\mathrm{Muk}}(K3))=8.
\]
Thus $\kappa+\kappa^!$ is not a universal number.  The Vol.~I
four-value ceiling is
$\{0,13,250/3,98/3\}$; the canonical cross-volume bucket is
$\{0,8,13,250/3,98/3\}$.
\end{proposition}

\begin{proof}
The formulas are scalar consequences of the Verdier-dual row, not of
the counit $\Omegach\Bbarch(\cA)\to\cA$.

For Heisenberg, $\kappa(\cH_k)=k$ and the Verdier partner has
$\kappa(\cH_k^!)=-k$, so the sum is zero.  For affine Kac--Moody,
\[
\kappa(V_k(\fg))=\frac{\dim(\fg)(k+h^\vee)}{2h^\vee},\qquad
k'=-k-2h^\vee,
\]
hence
\[
\kappa(V_{k'}(\fg))
=\frac{\dim(\fg)(-k-h^\vee)}{2h^\vee}
=-\kappa(V_k(\fg)).
\]
For the $\beta\gamma$ row,
\[
\kappa(\beta\gamma_\lambda)=6\lambda^2-6\lambda+1,\qquad
\kappa((\beta\gamma_\lambda)^!)
=-(6\lambda^2-6\lambda+1),
\]
again giving zero.  For Virasoro, the finite-type same-family
Verdier partner is $\mathrm{Vir}_{26-c}$ on the
$\mathsf M/\mathsf S$ surface, so
\[
\kappa(\mathrm{Vir}_c)+\kappa(\mathrm{Vir}_{26-c})
=\frac{c}{2}+\frac{26-c}{2}=13.
\]

For the principal $\mathcal W_3$ row,
$\kappa(\mathcal W_3)=c(H_3-1)=5c/6$ and the conductor is
$K_3=c+c'=4\cdot3^3-2\cdot3-2=100$; hence
$K^\kappa=(5/6)\cdot100=250/3$.  For the
Bershadsky--Polyakov row, $\kappa(\mathrm{BP}_k)=c(k)/6$ and
$c(k)+c(-k-6)=196$, giving $K^\kappa=196/6=98/3$.  For the
Mukai-enhanced K3 row, the Vol.~III input is
$c_+(\mathrm{Mukai}(K3))=4$, so the Mukai-doubling conductor gives
$K^\kappa=2c_+(\mathrm{Mukai}(K3))=8$. These are precisely the
values recorded in
Theorem~\ref{thm:census-witness-complementarity},
Proposition~\ref{prop:archetype-complementarity-bridge}, and
Convention~\ref{conv:theorem-c-bucket}.
\end{proof}

\begin{remark}[Three levels of Koszul pair]\label{rem:three-koszul-levels}
\index{Koszul pair!comparison of levels}
The construction above exhibits three levels of structure, each
strictly refining the previous:
\begin{enumerate}[leftmargin=2em]
\item \textbf{Classical} (Definition~\ref{def:koszul-pair-classical}).
 Input: graded algebras on a point. Bar object: $\barB(A) = T^c(s^{-1}\bar{A})$.
 Curvature: $m_0 = 0$. Main result: Priddy~\cite{Priddy70}.
\item \textbf{Chiral} (Definition~\ref{def:chiral-koszul-pair}).
 Input: chiral algebras on a curve~$X$. Bar object: $\barB_X(\cA)$
 on $\FM_n(X)$. Curvature: $\dzero^2 = 0$. Main result: Theorem~A
 (Theorem~\ref{thm:bar-cobar-isomorphism-main}).
\item \textbf{Modular} (Definition~\ref{def:modular-koszul-homotopy}).
 Input: chiral algebras + genus tower $g \geq 0$. Bar object:
 $\barB_g(\cA)$ on $\overline{\mathcal{M}}_g$. Scalar curvature:
 $\dfib^2 = \kappa(\cA)\cdot \omega_g$ on the uniform-weight lane,
 with completed CDG corrections in the multi-weight case. Main
 results: Theorems~B--D.
\end{enumerate}
Each level introduces new phenomena invisible to the previous one:
the chiral level brings integration kernels, OPE poles, and the
configuration-space geometry of $\FM_n(X)$ (Arnold relations,
logarithmic propagator); the modular level brings curvature,
central extensions, and anomalies. The classical and chiral levels
are \emph{not} identified by specializing to genus~$0$ on~$\bP^1$:
the chiral bar on~$\bP^1$ retains the full FM compactification
geometry, and the passage to the classical bar over a point requires
formal-disk restriction plus homotopy-transfer data (see the
discussion following Definition~\ref{def:chiral-koszul-pair}).
\end{remark}

\begin{proposition}[Three levels as MC at successive completions]
\label{prop:three-levels-mc-completion}
\ClaimStatusConditional
The three levels of Koszul structure (classical, chiral, and modular) correspond to the MC equation in successive completions of the convolution algebra:
\begin{enumerate}[label=\textup{(\roman*)}]
\item Classical: $\Theta_\cA^{(0,2)} \in \mathrm{MC}(\mathfrak{g}^{(0,\leq 2)}_\cA)$ (genus~$0$, degree~$2$).
\item Chiral: $\Theta_\cA^{(0,\bullet)} \in \mathrm{MC}(\mathfrak{g}^{(0)}_\cA)$ (genus~$0$, all degrees).
\item Modular: $\Theta_\cA \in \mathrm{MC}(\gAmod)$ (all genera, all degrees).
\end{enumerate}
The PBW filtration connects these levels: MC1 (PBW concentration) implies that the genus-$0$ MC element lifts to all genera via the strong completion-tower theorem (MC4, Theorem~\ref{thm:completed-bar-cobar-strong}).
\end{proposition}

\begin{proof}
At each level, the MC equation $D\Theta + \tfrac{1}{2}[\Theta,\Theta] = 0$ is the projection of $D_\cA^2 = 0$ (Theorem~\ref{thm:convolution-d-squared-zero}) to the corresponding truncation of the weight filtration on $\gAmod$.
Classical: retain only genus~$0$, degree~$2$ (quadratic).
Chiral: retain genus~$0$, all degrees (the full genus-$0$ shadow obstruction tower, with obstruction classes $o_{r+1}$ in the cyclic deformation complex).
Modular: the full convolution algebra with no truncation.
The PBW bridge (Theorem~\ref{thm:uniform-pbw-bridge}) reduces the modular MC equation to the chiral one on the associated graded, which the degree cutoff (Lemma~\ref{lem:degree-cutoff}) and the strong completion-tower theorem lift to all genera.
\end{proof}

\subsection{From classical to chiral: the PBW bridge}

Classical Koszul duality (Definition~\ref{def:koszul-pair-classical}, Chapter~\ref{ch:algebraic-foundations}) requires quadratic relations. Chiral algebras are almost never quadratic: the Virasoro OPE has a quartic pole, $W$-algebras have sixth-order poles, and Yangians have cubic relations. The PBW filtration (Chapter~\ref{chap:koszul-pairs}) resolves this by reducing the chiral problem to the classical one on the associated graded.

% ================================================================
% SECTION: CHIRAL HOCHSCHILD COHOMOLOGY FROM FIRST PRINCIPLES
% ================================================================

\section{Chiral Hochschild cohomology: construction from first principles}

\subsection{The chiral enveloping algebra}

\begin{definition}[Chiral enveloping algebra]
Let $\mathcal{A}$ be a unital associative chiral algebra on a smooth
curve~$X$, in a complete subcategory of right $\mathcal D_X$-modules
closed under the chiral tensor product.  The \emph{chiral enveloping
algebra} is
\[
\mathcal{A}^e = \mathcal{A} \boxtimes_{\mathcal{D}_X} \mathcal{A}^{\text{op}}
\]
where $\boxtimes_{\mathcal{D}_X}$ is the chiral tensor product over
the sheaf of differential operators and $\mathcal{A}^{\text{op}}$
has the opposite chiral multiplication
$Y^{\text{op}}(a,b)=(-1)^{|a||b|}Y(b,a)$ in the dg convention.
All tensor products below are completed in the PBW/admissible
topology when the bar filtration is infinite.
\end{definition}

\begin{definition}[Chiral bimodule]\label{def:chiral-bimodule}
An \emph{$\mathcal{A}$-bimodule} (in the chiral sense) is a $\mathcal{D}_X$-module $\mathcal{M}$ equipped with a left chiral action $\mu_L\colon j_*j^*(\mathcal{A} \boxtimes \mathcal{M}) \to \Delta_! \mathcal{M}$ and a right chiral action $\mu_R\colon j_*j^*(\mathcal{M} \boxtimes \mathcal{A}) \to \Delta_! \mathcal{M}$ satisfying the associativity and compatibility conditions:
\begin{enumerate}
\item Left associativity: $\mu_L \circ (\mathrm{id}_\mathcal{A} \boxtimes \mu_L) = \mu_L \circ (\mu_\mathcal{A} \boxtimes \mathrm{id}_\mathcal{M})$ (as chiral operations on $\mathcal{A} \boxtimes \mathcal{A} \boxtimes \mathcal{M}$);
\item Right associativity (analogously);
\item Left-right compatibility: $\mu_L \circ (\mathrm{id}_\mathcal{A} \boxtimes \mu_R) = \mu_R \circ (\mu_L \boxtimes \mathrm{id}_\mathcal{A})$.
\end{enumerate}
Equivalently, $\mathcal{M}$ is a module for the chiral enveloping algebra $\mathcal{A}^e$.
\end{definition}

\begin{lemma}[Well-definedness of the chiral enveloping algebra;
\ClaimStatusProvedHere]
\label{lem:chiral-enveloping-well-defined}
Under the closure and completion hypotheses in the preceding
definition, $\mathcal{A}^e$ carries a natural associative chiral
algebra structure, and chiral $\mathcal{A}^e$-modules are the same
as chiral $\mathcal{A}$-bimodules.
\end{lemma}

\begin{proof}
The tensor product exists by the closure hypothesis; no projective
resolution is being smuggled into the definition.  On disjoint
configurations the product is
\[
Y^e((a_1 \otimes a_2), (b_1 \otimes b_2))(z)
=(-1)^{|a_2||b_1|}Y(a_1,b_1)(z)\otimes
Y^{\mathrm{op}}(a_2,b_2)(z),
\]
with the displayed Koszul sign coming from moving the second tensor
factor past the first.  Associativity is the tensor product of the
associativity diagrams for $\cA$ and $\cA^{\mathrm{op}}$; locality is
preserved because the pole order is bounded by the sum of the pole
orders in the two factors.  Unpacking a left action of this product
gives commuting left and right chiral $\cA$-actions, and the reverse
construction tensors the two actions. \qedhere
\end{proof}

\subsection{The bar resolution for chiral algebras}

The chiral Hochschild complex is obtained from the relative
two-sided bar resolution of $\mathcal{A}$ as an
$\mathcal{A}^e$-module.  This is a resolution in the completed
derived category; it is not the Koszul bar coalgebra
$\barB^{\mathrm{ch}}(\cA)$ and it does not identify
$\cA^!$.

\begin{definition}[Two-sided chiral Hochschild bar resolution]
The \emph{two-sided chiral Hochschild bar resolution} of
$\mathcal{A}$ is
\[
\cdots \to \mathcal{A}^{\boxtimes 4} \xrightarrow{d_3} \mathcal{A}^{\boxtimes 3} \xrightarrow{d_2} \mathcal{A}^{\boxtimes 2} \xrightarrow{d_1} \mathcal{A} \to 0
\]
where the outer two factors carry the $\mathcal A^e$-action and
\[
d_n(a_0 \otimes \cdots \otimes a_{n+1}) = \sum_{i=0}^n (-1)^i a_0 \otimes \cdots \otimes Y(a_i, a_{i+1}) \otimes \cdots \otimes a_{n+1}
\]
with Koszul signs inserted in the dg convention.  The normalized
complex is used when degeneracies are split off.
\end{definition}

\begin{theorem}[Relative exactness of the two-sided chiral bar
resolution; \ClaimStatusProvedHere]
\label{thm:chiral-bar-resolution-exact}
\textup{[Regime: completed unital chiral algebra
\textup{(}Convention~\textup{\ref{conv:regime-tags})}.]}

Let $\mathcal A$ be unital, augmented, and admissibly complete, and
assume the chiral tensor product is exact on the bar-filtration
graded pieces.  The normalized two-sided chiral bar object is a
relative semi-free $\mathcal A^e$-resolution of $\mathcal A$ in the
completed derived category.  After forgetting the
$\mathcal A^e$-action the augmented complex is contractible.
This exactness is the Hochschild-resolution input; Koszulness is the
separate assertion that the one-sided conilpotent bar coalgebra
$\barB^{\mathrm{ch}}(\cA)$ resolves the Koszul-dual coalgebra
$\cA^i$.
\end{theorem}

\begin{proof}
The simplicial two-sided bar object has face maps given by the
chiral product and degeneracies given by insertion of the unit.
The extra degeneracy
\[
h_n(a_0 \otimes \cdots \otimes a_n) = 1 \otimes a_0 \otimes \cdots \otimes a_n
\]
gives the standard contraction of the augmented simplicial object.
On homogeneous tensors the only non-cancelling face is the one
multiplying the inserted unit:
\begin{align}
(d_{n+1} \circ h_n)(a_0 \otimes \cdots \otimes a_n) &= d_{n+1}(1 \otimes a_0 \otimes \cdots \otimes a_n) \\
&= a_0 \otimes \cdots \otimes a_n + \sum_{i=0}^{n-1} (-1)^{i+1} 1 \otimes a_0 \otimes \cdots \otimes Y(a_i, a_{i+1}) \otimes \cdots
\end{align}
while
\[
(h_{n-1} \circ d_n)(a_0 \otimes \cdots \otimes a_n) = -\sum_{i=0}^{n-1} (-1)^{i+1} 1 \otimes a_0 \otimes \cdots \otimes Y(a_i, a_{i+1}) \otimes \cdots
\]
up to the Koszul signs forced by the suspended convention.  Thus
$d h+h d=\mathrm{id}$ after forgetting the outer action.  The
normalized bar object is built from induced
$\mathcal A^e$-modules, hence is relative semi-free. Exactness of
the completed tensor product on associated graded pieces lifts the
contraction from finite bar degree to the completed object. \qedhere
\end{proof}

\subsection{Definition and computation of chiral Hochschild cohomology}

\begin{definition}[Chiral Hochschild cohomology, algebraic Ext form]
\label{def:chirHoch-Ext}
\index{Hochschild cohomology!chiral!algebraic Ext form}

In the completed chiral bimodule category, the \emph{chiral
Hochschild cohomology} of $\mathcal{A}$ with coefficients in an
$\mathcal{A}$-bimodule $M$ is
\[
\ChirHoch^n(\mathcal{A}, M)
= \operatorname{Ext}^n_{\mathcal{A}^e}(\mathcal{A}, M).
\]
When $M = \mathcal{A}$, we write simply $\ChirHoch^n(\mathcal{A})$.
This is the \emph{algebraic Ext model} of the chiral Hochschild
complex; the geometric Fulton--MacPherson model
$C^\bullet_{\mathrm{chiral}}(\cA)$
\textup{(}Definition~\ref{def:chiral-hochschild-complex}, in
Chapter~\textup{\ref{chap:deformation-theory}}\textup{)} computes
the same chiral derived endomorphism centre
\[
 Z^{\mathrm{der}}_{\mathrm{ch}}(\cA)
 = R\!\operatorname{Hom}^{\mathrm{ch}}_{\cA^e}(\cA,\cA)
\]
under
Theorem~\textup{\ref{thm:geometric-chiral-hochschild}}. Per
Remark~\textup{\ref{rem:three-hochschild-discipline}} in
Chapter~\textup{\ref{chap:deformation-theory}}, the chiral
$\ChirHoch^*$ is distinct from the topological Hochschild
$\mathrm{HH}^*_{\mathrm{top}}$ of an $E_1$-algebra in spectra
\textup{(}Deligne $E_2$\textup{)} and from the categorical Hochschild
$\mathrm{HH}^*_{\mathrm{cat}}$ of a dg category \textup{(}CY shifted
Poisson\textup{)}. The geometry $X$ selects $\ChirHoch^*$.
\end{definition}

To compute this explicitly, we apply $\text{Hom}_{\mathcal{A}^e}(-, M)$ to the bar resolution:

\begin{theorem}[Chiral Hochschild complex; \ClaimStatusProvedHere]
\label{thm:chiral-hochschild-complex}
For every complete chiral $\mathcal A$-bimodule~$M$ satisfying the
same admissibility hypotheses, the derived endomorphism complex
$R\!\operatorname{Hom}_{\mathcal A^e}^{\mathrm{ch}}(\mathcal A,M)$
is computed by the normalized chiral Hochschild cochain complex
\[
0 \to \text{Hom}_{\mathcal{D}_X}(\mathcal{A}, M) \xrightarrow{\delta_0} \text{Hom}_{\mathcal{D}_X}(\mathcal{A}^{\otimes 2}, M) \xrightarrow{\delta_1} \text{Hom}_{\mathcal{D}_X}(\mathcal{A}^{\otimes 3}, M) \to \cdots
\]
where the differential $\delta_n$ is:
\begin{align}
(\delta_n f)(a_0, \ldots, a_{n+1}) = &Y(a_0, f(a_1, \ldots, a_{n+1})) \\
&+ \sum_{i=1}^n (-1)^i f(a_0, \ldots, Y(a_i, a_{i+1}), \ldots, a_{n+1}) \\
&+ (-1)^{n+1} Y(f(a_0, \ldots, a_n), a_{n+1})
\end{align}
\end{theorem}
\begin{proof}
Apply $R\!\mathrm{Hom}_{\mathcal A^e}^{\mathrm{ch}}(-,M)$ to the
relative semi-free resolution of
Theorem~\ref{thm:chiral-bar-resolution-exact}.  The outer
$\mathcal A$-actions produce the first and last terms in
$\delta_n$, and the internal bar faces produce the middle sum.
\end{proof}

\subsection{Geometric realization via configuration spaces}

\begin{theorem}[Geometric model of chiral Hochschild cohomology; \ClaimStatusProvedHere]
\label{thm:geometric-chiral-hochschild}
Assume, in addition, that the chiral operations of~$\mathcal A$
extend logarithmically to the Fulton--MacPherson compactifications
and that the chosen finite-type stratum is acyclic for the
\v{C}ech-to-global functor on the sheaves below.  Then the algebraic
Ext model is filtered quasi-isomorphic to the total
Fulton--MacPherson logarithmic complex with degree-$n$ term
\[
\Gamma\left(\overline{C}_{n+1}(X),
\mathcal{H}om_{\mathcal{D}_X}
(\mathcal{A}^{\boxtimes n+1},\mathcal{A})
\otimes \Omega^n_{\log}\right),
\]
and cohomology equal to $\ChirHoch^n(\cA)$.  This comparison computes
$Z^{\mathrm{der}}_{\mathrm{ch}}(\cA)$; it does not produce
$\cA^!$.
\end{theorem}

\begin{proof}
The comparison is the following chain of identifications.

\emph{Step 1}: An $\mathcal{A}^e$-linear map
$f: \mathcal{A}^{\otimes n+1} \to \mathcal{A}$ is the same thing as
a chiral multilinear operation compatible with the two outer
$\mathcal A$-actions:
\[
f(Y(a, z)b_1, \ldots, b_n) = Y(a, z)f(b_1, \ldots, b_n)
\]
\[
f(b_1, \ldots, b_n, Y(b, w)) = Y(f(b_1, \ldots, b_n), w)b
\]

\emph{Step 2}: On the open configuration space $C_{n+1}(X)$ these
conditions identify a cochain with a section of the stated internal
Hom sheaf.

\emph{Step 3}: The logarithmic extension hypothesis extends precisely
those sections to $\overline{C}_{n+1}(X)$ with logarithmic
singularities along the boundary divisors.  No finite-type dual of
the bar coalgebra is used in this extension.

\emph{Step 4}: The chiral Hochschild differential
$\delta_n f(a_0, \ldots, a_{n+1})$ involves composing~$f$ with the
chiral product at adjacent points. On $\overline{C}_{n+2}(X)$, each
term in $\delta_n$ corresponds to the residue of $f$ along a boundary
divisor $D_{ij} = \{z_i = z_j\}$ (where two adjacent points collide).
The de~Rham differential on $\Omega^n_{\log}$ restricts to these
boundary strata by Stokes' theorem, and the residue along
$D_{ij}$ extracts the OPE coefficient: $\mathrm{Res}_{z_i = z_j}
[f \cdot \omega_{\log}] = Y(a_i, a_j) \cdot f|_{D_{ij}}$.
The alternating signs in $\delta_n$ correspond to the orientations
of the boundary divisors in $\partial\overline{C}_{n+2}(X)$
(cf.\ the bar nilpotency argument,
Theorem~\ref{thm:bar-nilpotency-complete}).  The acyclicity
assumption for the chosen finite-type stratum lets global sections
preserve the resulting filtered quasi-isomorphism. \qedhere
\end{proof}

% ================================================================
% SECTION: THE CHIRAL GERSTENHABER STRUCTURE
% ================================================================

\section{The chiral Gerstenhaber structure}

Classical topological Hochschild cohomology carries a Gerstenhaber algebra structure~\cite{Ger63}: a graded commutative cup product and a degree~$-1$ Lie bracket, compatible via the Leibniz rule. Both extend to the chiral setting.

\subsection{The cup product}

\begin{definition}[Cup product on chiral Hochschild cohomology]
For $f \in \ChirHoch^p(\mathcal{A})$ and $g \in \ChirHoch^q(\mathcal{A})$, define their cup product:
\[
(f \cup g)(a_0, \ldots, a_{p+q}) = \text{Res}_{z_p \to w_0} f(a_0, \ldots, a_p)(z_0, \ldots, z_p) \cdot g(a_{p+1}, \ldots, a_{p+q})(w_0, \ldots, w_{q-1})
\]
where the residue is taken as the $p$-th point approaches the position of the $(p+1)$-st point.
\end{definition}

\begin{proposition}[Properties of cup product; \ClaimStatusProvedHere]
\label{prop:cup-product-properties}
On the Hochschild cohomology groups computed by
Theorem~\textup{\ref{thm:chiral-hochschild-complex}}, the cup
product satisfies:
\begin{enumerate}
\item \emph{Associativity}: $(f \cup g) \cup h = f \cup (g \cup h)$
\item \emph{Graded commutativity}: $f \cup g = (-1)^{|f||g|} g \cup f$
\item \emph{Unit}: The identity element $1 \in \ChirHoch^0(\mathcal{A})$ is a unit for $\cup$
\end{enumerate}
\end{proposition}

\begin{proof}
\emph{Associativity}: Both $(f \cup g) \cup h$ and $f \cup (g \cup h)$ involve taking residues at collision points. The order of residues does not matter by the residue theorem on $\overline{C}_{n}(X)$.

\emph{Graded commutativity}: This is a cohomological statement.  On
cochains the two products are related by the Eilenberg--Zilber
homotopy coming from the transposition of the two point blocks in
$\overline{C}_n(X)$; the Koszul sign is the sign of the corresponding
permutation of logarithmic forms.

\emph{Unit}: The identity in $\ChirHoch^0$ is the identity map $\mathcal{A} \to \mathcal{A}$, which acts trivially under cup product. \qedhere
\end{proof}

\subsection{The chiral Lie bracket}

The chiral Lie bracket uses residue composition:

\begin{definition}[Chiral Lie bracket]
For $f \in \ChirHoch^p(\mathcal{A})$ and $g \in \ChirHoch^q(\mathcal{A})$, define:
\[
\{f, g\}_c = f \circ_c g - (-1)^{(p-1)(q-1)} g \circ_c f
\]
where the chiral composition $\circ_c$ is:
\begin{multline}
(f \circ_c g)(a_0, \ldots, a_{p+q-1})
= \sum_{i=0}^{p-1} (-1)^{i(q-1)}
\operatorname{Res}_{w \to z_i}\\
f\bigl(a_0, \ldots, a_i,\, g(a_{i+1}, \ldots, a_{i+q})(w),\,
a_{i+q+1}, \ldots\bigr)
\end{multline}
\end{definition}

\begin{theorem}[Chiral Gerstenhaber algebra {\cite{Ger63, Tamarkin00}}; \ClaimStatusProvedElsewhere]
\label{thm:chiral-gerstenhaber-kps}
Under the same conilpotent/completed bar hypotheses, the cohomology
$\ChirHoch^*(\mathcal{A})$ with operations $(\cup, \{-,-\}_c)$ forms
a Gerstenhaber algebra:
\begin{enumerate}
\item \emph{Chiral Jacobi identity}:
\[
\{f, \{g, h\}_c\}_c = \{\{f, g\}_c, h\}_c + (-1)^{(|f|-1)(|g|-1)} \{g, \{f, h\}_c\}_c
\]
\item \emph{Chiral Leibniz rule}:
\[
\{f, g \cup h\}_c = \{f, g\}_c \cup h + (-1)^{(|f|-1)|g|} g \cup \{f, h\}_c
\]
\end{enumerate}
This is the chiral analogue of Gerstenhaber's classical theorem~\cite{Ger63}.
The bracket is the commutator on the dg Lie algebra of continuous
coderivations of the conilpotent bar coalgebra
$\bar{B}^{\mathrm{ch}}(\mathcal{A})$; the Hochschild complex maps to
that coderivation complex by the universal twisting morphism.  The
Leibniz rule is the compatibility of this commutator with the
composition product on self-extensions of the diagonal bicomodule
(Tamarkin~\cite{Tamarkin00}; cf.\ the classical statement in
\cite[Thm.~9.3.2]{LV12}).
\end{theorem}

% ================================================================
% SECTION: HIGHER STRUCTURES - A-INFINITY AND L-INFINITY
% ================================================================

\section{\texorpdfstring{Higher structures: $A_\infty$ and $L_\infty$ on chiral Hochschild cohomology}{Higher structures: A-infinity and L-infinity on chiral Hochschild cohomology}}

\subsection{\texorpdfstring{The $A_\infty$ structure}{The A-infinity structure}}

\begin{theorem}[\texorpdfstring{$A_\infty$}{A-infinity} operations from the chiral brace model; \ClaimStatusProvedHere]
\label{thm:ainfty-chiral-hochschild}
\leavevmode\\
Choose a cellular chain model for the little-discs operad inside the
Fulton--MacPherson logarithmic compactifications used in
Theorem~\textup{\ref{thm:geometric-chiral-hochschild}}.  The
normalized chiral Hochschild cochain complex carries brace
operations; restriction to the associative cellular suboperad gives
operations
\[
m_n: \ChirHoch^{i_1} \otimes \cdots \otimes \ChirHoch^{i_n} \to \ChirHoch^{i_1 + \cdots + i_n + 2 - n}
\]
satisfying the $A_\infty$ relations (Stasheff--Loday--Vallette convention):
\[
\sum_{r+s+t=n} (-1)^{rs+t}\, m_{r+1+t}(\mathrm{id}^{\otimes r} \otimes m_s \otimes \mathrm{id}^{\otimes t}) = 0
\]
where the sign $(-1)^{rs+t}$ arises from the Koszul sign rule on
suspended generators.  The construction is a chain-level model for
the derived centre; after passing to cohomology it recovers the
Gerstenhaber operations of
Theorem~\textup{\ref{thm:chiral-gerstenhaber-kps}}.
\end{theorem}

\begin{proof}[Construction of higher operations]
The brace operation indexed by a rooted planar tree inserts
cochains at the corresponding ordered collision strata and composes
with the chiral product through the universal twisting morphism.
For a top-dimensional cell~$\sigma$ in the chosen associative
cellular suboperad, define $m_\sigma$ by summing the residue
operations over the boundary divisors met by~$\sigma$.  The operation
$m_n$ is the sum over the $n$-input cells with the orientation signs
of the cellular chain model.

The $A_\infty$ relation is the cellular boundary identity.  The
codimension-one boundary of a compactified configuration cell is a
union of products of smaller cells,
\[
\partial \overline{C}_n(\mathbb{P}^1) = \bigcup_{i+j=n+1} \bigcup_{I \sqcup J = [n]} \overline{C}_i(\mathbb{P}^1) \times \overline{C}_j(\mathbb{P}^1)
\]
and each boundary component contributes the corresponding insertion
$m_i\circ_k m_j$.  The sign $(-1)^{rs+t}$ is the
Stasheff--Loday--Vallette sign obtained after desuspending the
inputs (see~\cite{LV12}, \S9.2.8).  Since the oriented cellular
boundary squares to zero, the sum of all insertions is zero. \qedhere
\end{proof}

\subsection{\texorpdfstring{The $L_\infty$ structure}{The L-infinity structure}}

By the commutator construction (antisymmetrizing the $A_\infty$ operations via the natural operad map $\mathrm{Lie} \to \mathrm{Ass}$), an $A_\infty$ structure induces an $L_\infty$ structure:

\begin{theorem}[\texorpdfstring{$L_\infty$}{L-infinity} structure {\cite{LV12}}; \ClaimStatusProvedElsewhere]
\label{thm:linfty-chiral-hochschild}
The shifted complex $\ChirHoch^{*-1}(\mathcal{A})[-1]$ carries an $L_\infty$ structure with brackets:
\[
\ell_n: \Lambda^n \ChirHoch^{*-1} \to \ChirHoch^{*-1}[2-n]
\]
related to the $A_\infty$ operations by:
\[
\ell_n(f_1, \ldots, f_n) = \sum_{\sigma \in S_n} \frac{\text{sign}(\sigma)}{n!} m_n(f_{\sigma(1)}, \ldots, f_{\sigma(n)})
\]
\end{theorem}

% ================================================================
% SECTION: PERIODICITY PHENOMENA
% ================================================================

\section{Periodicity in chiral Hochschild cohomology}

\subsection{Statement and first properties}

\begin{theorem}[Virasoro amplitude-two Hochschild class; \ClaimStatusConditional]\label{thm:periodicity-virasoro}
\label{thm:virasoro-periodicity}
\index{periodicity!Virasoro}
\index{Virasoro algebra!periodicity}
For the Virasoro algebra $\mathrm{Vir}_c$ at generic central charge $c$
(avoiding the minimal model values
$c_{p,q} = 1 - 6(p-q)^2/(pq)$ and $c = 1$), assume the following
three inputs on the same completed curve-level Hochschild surface:
\begin{enumerate}[label=\textup{(V\arabic*)}]
\item the Kac-determinant genericity input kills the higher internal
Ext terms in the Virasoro block used by the normalized chiral bar
filtration;
\item the Beilinson--Drinfeld/Fulton--MacPherson comparison identifies
the degree-$\leq2$ part of the chiral Hochschild complex with the
corresponding relative continuous Virasoro model;
\item the de~Rham amplitude bound of
Theorem~\textup{\ref{thm:hochschild-polynomial-growth}} applies to
this surface, so no curve-level chiral Hochschild class survives in
degree $>2$.
\end{enumerate}
Then there exists a class $\Theta \in \ChirHoch^2(\mathrm{Vir}_c)$
such that cup product with $\Theta$ induces the single nonzero
amplitude-two isomorphism
\[
\ChirHoch^0(\mathrm{Vir}_c) \xrightarrow{\;\cup\,\Theta\;}
\ChirHoch^2(\mathrm{Vir}_c).
\]
Moreover $\ChirHoch^1(\mathrm{Vir}_c)=0$ and
$\ChirHoch^n(\mathrm{Vir}_c)=0$ for $n>2$ on this curve-level
chiral Hochschild surface.  The unbounded polynomial
Gel'fand--Fuchs periodicity belongs to continuous Lie algebra
cohomology before the chiral de~Rham amplitude truncation.
\end{theorem}

\begin{proof}
The compactification $\overline{\mathcal{M}}_{0,4} \cong \mathbb{P}^1$
carries the Weil--Petersson class
\index{Weil--Petersson class}
$\Theta = [\omega_{WP}] \in
H^2(\overline{\mathcal{M}}_{0,4}) \cong \mathbb{C}$, the fundamental
class of the unique boundary divisor.  Input~(V2) maps this class to
a Hochschild cocycle because the logarithmic de~Rham differential and
the chiral Hochschild differential correspond under residues along
the boundary divisors.  Thus $\Theta$ defines a class in
$\ChirHoch^2(\mathrm{Vir}_c)$.

The bar filtration gives a spectral sequence
\[
E_2^{p,q} = H^p(\overline{\mathcal{M}}_{0,\bullet}) \otimes
H^q(\mathrm{Vir}_c\text{-mod})
\;\Longrightarrow\; \ChirHoch^{p+q}(\mathrm{Vir}_c).
\]
Input~(V1) says that the $q>0$ line vanishes on the relevant generic
block; hence the curve-level contribution is concentrated in the
$q=0$ line.  Input~(V3) kills all terms of total degree $>2$.
The remaining degree-zero class is the scalar centre, the
degree-one term vanishes by (V1), and the degree-two class is the
image of $\Theta$.

\index{Gel'fand--Fuchs theorem}
The Gel'fand--Fuchs theorem
\cite[Ch.~2, \S2]{Fuks86} computes the continuous relative
cohomology of the Lie algebra
$L_1 = t\,\mathbb{C}[[t]]\,d/dt$ of formal vector fields vanishing at
the origin:
\[
H^*_{\mathrm{cont}}(L_1;\, \mathbb{C}) = \mathbb{C}[\Theta],
\qquad \Theta \in H^2.
\]
Input~(V2) identifies the degree-two generator with the
Weil--Petersson boundary class in the geometric bar model; input~(V3)
truncates the polynomial continuous cohomology to the curve-level
amplitude range.  Therefore
\[
\ChirHoch^0 = \mathbb{C},
\qquad
\ChirHoch^1 = 0,
\qquad
\ChirHoch^2 = \mathbb{C}\cdot\Theta.
\]
The Gel'fand--Fuchs identification
$H^*_{\mathrm{cont}}(L_1;\,\mathbb{C}) = \mathbb{C}[\Theta]$
describes the continuous Lie algebra cohomology (unbounded); the
chiral Hochschild cohomology on a curve truncates to the
$\{0,1,2\}$~range.

The cup product
$\cup\,\Theta\colon \ChirHoch^0 \to \ChirHoch^2$ maps $1 \mapsto \Theta$,
which is an isomorphism $\mathbb{C} \xrightarrow{\sim} \mathbb{C}$.
The exclusion of $c=1$ and of the minimal-model values is precisely
the genericity assumption needed for (V1).
\end{proof}

\begin{remark}[Scope]
The open mathematical condition is not the Gel'fand--Fuchs calculation itself;
that calculation is the unbounded continuous Lie algebra input.  The
load-bearing comparison is the passage from that local continuous
model to the curve-level chiral Hochschild complex, together with the
amplitude truncation of
Theorem~\ref{thm:hochschild-polynomial-growth}.  The conclusion is an
amplitude-two statement on $Z^{\mathrm{der}}_{\mathrm{ch}}
(\mathrm{Vir}_c)$, not a strict all-degree periodicity theorem.
\end{remark}

\subsection{Periodicity for other chiral algebras}

\begin{theorem}[Critical formal-disk Hochschild model; \ClaimStatusConditional]
\label{thm:affine-periodicity-critical}
\index{periodicity!affine Kac--Moody}
Let $\mathfrak{g}$ be a finite-dimensional simple Lie algebra of rank~$r$
with exponents $m_1 \leq \cdots \leq m_r$.
At critical level $k = -h^{\vee}$, write
$\ChirHoch_{\mathrm{loc}}^*$ for the formal-disk continuous Lie
algebra model obtained before imposing the global curve-level
de~Rham amplitude bound.  Assume the Beilinson--Drinfeld
formal-disk comparison identifies this local model with
continuous cohomology of $\mathfrak g\otimes t\bC[[t]]$.  Then
\[
\ChirHoch_{\mathrm{loc}}^*(\hat{\mathfrak{g}}_{-h^{\vee}}) \;\cong\;
\Lambda^*(P_1, \ldots, P_r) \otimes \bC[\Theta_1, \ldots, \Theta_r]
\]
where $\deg P_i = 2m_i + 1$ (odd generators) and $\deg\Theta_i = 2(m_i + 1)$
(even generators). In particular:
\begin{enumerate}
\item For $\mathfrak{g} = \mathfrak{sl}_2$ (rank $1$, exponent $m_1 = 1$):
$\ChirHoch_{\mathrm{loc}}^* = \Lambda(P_1) \otimes \bC[\Theta_1]$
with $\deg P_1 = 3$, $\deg\Theta_1 = 4$, and
$\ChirHoch_{\mathrm{loc}}^{n+4} \cong
\ChirHoch_{\mathrm{loc}}^n$ for all $n \geq 0$.
\item For rank $r > 1$: the Poincar\'e series is
$\prod_{i=1}^r (1 + t^{2m_i+1})/(1-t^{2(m_i+1)})$, giving polynomial
growth $\dim \ChirHoch_{\mathrm{loc}}^n = O(n^{r-1})$ but
\emph{not} strict periodicity when the exponents are not all equal.
\end{enumerate}
This is the local continuous model.  The global curve-level chiral
Hochschild object $Z^{\mathrm{der}}_{\mathrm{ch}}$ is obtained only
after the logarithmic/Fulton--MacPherson comparison and its amplitude
truncation are imposed.  The statement does not identify the
critical affine algebra with its Koszul dual.
\end{theorem}

\begin{proof}
\emph{Step 1: critical local comparison.}
At critical level $k = -h^\vee$, the bar complex is \emph{uncurved}:
the curvature $m_0 \in \bar{B}^0$ has coefficient
$\kappa = (k+h^\vee)\dim(\fg)/(2h^\vee)$, hence vanishes when
$k=-h^\vee$.  This is not a commutativity statement about the affine VOA:
the critical current OPE still contains the simple-pole bracket term and
the level double pole $k(\cdot,\cdot)/(z-w)^2$.  What becomes commutative
is the Feigin--Frenkel centre
$\mathfrak{z}(\widehat{\fg}_{-h^\vee})\cong
\mathrm{Fun}\,\mathrm{Op}_{\fg^\vee}(D)$.  The Sugawara construction is
undefined at critical level, so no Sugawara-normalised stress tensor is used
to prove curvature vanishing.  The level-shifting duality
(Corollary~\ref{cor:level-shifting-part1}) confirms that the dual level
$k' = -k - 2h^\vee = -h^\vee$ is the unique fixed point of the
Feigin--Frenkel involution.

By the formal-disk comparison hypothesis, together with the
Feigin--Frenkel identification of the critical chiral centre
\cite{Feigin-Frenkel},
\[
\ChirHoch_{\mathrm{loc}}^*(\hat{\mathfrak{g}}_{-h^\vee}) \;\cong\;
H^*_{\mathrm{Lie,cont}}(\mathfrak{g} \otimes t\bC[[t]];\, \bC).
\]

\emph{Step 2: Gel'fand--Fuchs computation.}
The continuous Lie algebra cohomology of $\mathfrak{g} \otimes t\bC[[t]]$
is computed by the Fuks--Feigin--Tsygan theorem (\cite{FT87}; see also
Fuks \cite{Fuks86}, Ch.~3):
\[
H^*_{\mathrm{Lie,cont}}(\mathfrak{g} \otimes t\bC[[t]];\, \bC) \;\cong\;
\Lambda^*(P_1, \ldots, P_r) \otimes \bC[\Theta_1, \ldots, \Theta_r]
\]
where $\deg P_i = 2m_i + 1$ and $\deg \Theta_i = 2(m_i + 1)$.
The odd generators $P_i$ correspond to primitive elements of
$H^*(\mathfrak{g}; \bC)$ (Chevalley--Koszul), shifted by the loop grading;
the even generators $\Theta_i$ arise from the $t$-adic filtration and
correspond to Chern--Weil classes.

\emph{Step 3: Periodicity for rank~$1$.}
For $\mathfrak{sl}_2$ ($r = 1$, $m_1 = 1$):
$\ChirHoch_{\mathrm{loc}}^* = \Lambda(P_1) \otimes \bC[\Theta_1]$
with Poincar\'e series $(1 + t^3)/(1 - t^4)$. Cup product with
$\Theta_1 \in \ChirHoch_{\mathrm{loc}}^4$ gives the isomorphism
$\ChirHoch_{\mathrm{loc}}^n \xrightarrow{\sim}
\ChirHoch_{\mathrm{loc}}^{n+4}$
for all $n \geq 0$: injectivity holds because $\bC[\Theta_1]$ is a
polynomial ring (no relations), and surjectivity because every element of
$\ChirHoch_{\mathrm{loc}}^{n+4}$ is divisible by $\Theta_1$ modulo
the exterior generator~$P_1$.

For rank $r > 1$, the Poincar\'e series
$\prod_{i=1}^r (1 + t^{2m_i+1})/(1-t^{2(m_i+1)})$ grows
polynomially of degree~$r-1$, so strict periodicity fails when
$r > 1$ (verified explicitly for $\mathfrak{sl}_3$ in
Remark~\ref{rem:affine-periodicity-analysis}).
\end{proof}

\begin{remark}[Status and rank distinction]\label{rem:rank-distinction}
\label{rem:affine-periodicity-analysis}
\emph{Rank~$1$.}
For $\hat{\mathfrak{sl}}_2$ at $k = -2$,
$\ChirHoch_{\mathrm{loc}}^* =
\Lambda(P_1) \otimes \bC[\Theta_1]$ with $\deg P_1 = 3$,
$\deg\Theta_1 = 4$, giving local $4$-periodicity
$\ChirHoch_{\mathrm{loc}}^{n+4}\cong
\ChirHoch_{\mathrm{loc}}^n$.  The cup product with $\Theta_1$ is the
local periodicity operator; the Virasoro theorem above is instead an
amplitude-two curve-level statement.  The centre
$\mathfrak{z}(\hat{\mathfrak{sl}}_{2,-2}) \cong
\mathrm{Fun}\,\mathrm{Op}_{\mathfrak{sl}_2}$ gives the geometric
origin of~$\Theta_1$.

\emph{Rank~$r > 1$.}
The local cohomology ring has $r$ distinct polynomial degrees
$2(m_i + 1)$; when not all equal, strict periodicity fails. For
$\hat{\mathfrak{sl}}_3$ at $k = -3$: $\deg\Theta_1 = 4$,
$\deg\Theta_2 = 6$, and
$\dim \ChirHoch_{\mathrm{loc}}^{12}=3\neq1=
\dim \ChirHoch_{\mathrm{loc}}^6$, violating $6$-periodicity. The
correct local statement is polynomial growth
$\dim \ChirHoch_{\mathrm{loc}}^n=O(n^{r-1})$.
\end{remark}

\begin{remark}[Provenance]\label{rem:periodicity-provenance}
The proof assembles three published results:
(i)~the Beilinson--Drinfeld comparison between the formal-disk
chiral Hochschild model and continuous Lie algebra cohomology
(\ClaimStatusConditional,
\cite{BD04}, Thm.~4.5.2);
(ii)~the Gel'fand--Fuchs/Feigin--Tsygan computation of
$H^*_{\mathrm{cont}}(\mathfrak{g} \otimes t\bC[[t]]; \bC)$
(\ClaimStatusProvedElsewhere, \cite{FT87});
(iii)~the uncurved bar complex structure at critical level, which
follows from the critical-level fixed-point property
(Corollary~\ref{cor:level-shifting-part1}).
These inputs give the stated critical-level local periodicity or
polynomial-growth profile for every simple~$\mathfrak{g}$.
\end{remark}

\begin{proposition}[Chiral Hochschild profile transport for same-type finite-type pairs; \ClaimStatusConditional]
\label{prop:periodicity-same-type}
\index{periodicity!preservation under Koszul duality}
For a finite-type Verdier-Koszul pair
$(\mathcal{A}, \mathcal{A}^!)$ where both algebras have the same type
\textup{(}both Kac--Moody, both Virasoro, or both
$\mathcal{W}_N$\textup{)}, assume both sides lie on the same
geometric Hochschild comparison surface, the curve-level amplitude
bound applies to both, and the finite-type derived equivalence is
Morita-compatible with that comparison.  For the $\mathcal W_N$ row,
also assume the relevant DS/bar transport identifying the same-family
Verdier partner.  Then the chiral Hochschild cohomology rings
$\ChirHoch^*(\mathcal{A})$ and $\ChirHoch^*(\mathcal{A}^!)$ have the
same growth type:
\begin{enumerate}[label=\textup{(\roman*)}]
\item \emph{Virasoro:} $\ChirHoch^*(\mathrm{Vir}_c)$ and $\ChirHoch^*(\mathrm{Vir}_{26-c})$ are both concentrated in $\{0, 2\}$ with $P(t) = 1 + t^2$ \textup{(}Theorem~\textup{\ref{thm:virasoro-hochschild}}\textup{)}.
\item \emph{Kac--Moody:} $\ChirHoch^*(\widehat{\mathfrak{g}}_k)$ and $\ChirHoch^*(\widehat{\mathfrak{g}}_{k'})$ are both concentrated in $\{0, 1, 2\}$ with $P(t) = 1 + \dim(\mathfrak{g})\,t + t^2$ at generic level \textup{(}the Hilbert polynomial depends on~$\mathfrak{g}$, not on~$k$\textup{)}.
\item \emph{$\mathcal{W}_N$:} $\ChirHoch^*(\mathcal{W}_N^k)$ and $\ChirHoch^*(\mathcal{W}_N^{k'})$ are both concentrated in $\{0, 1, 2\}$ with the same Hilbert polynomial, since the conformal weights $2, 3, \ldots, N$ are level-independent.
\end{enumerate}
\end{proposition}

\begin{proof}
In each case, the local continuous model records the root data of
$\mathfrak{g}$ and the conformal weights of the generators, neither
of which changes under the Feigin--Frenkel involution
$k \mapsto k' = -k - 2h^\vee$. The assumed global chiral Hochschild
comparison imposes the same curve-level amplitude bound on both
members of the same-type pair. Part~(i) is a direct consequence of
Theorem~\ref{thm:virasoro-periodicity}; parts~(ii) and~(iii) use the
level-independence of the exponents and strong-generator weights,
together with the explicit same-surface assumptions in the statement.
\end{proof}

\begin{corollary}[chiral Hochschild transport under finite-type Koszul duality; \ClaimStatusConditional]
\label{cor:hochschild-ring-koszul}
\index{Hochschild cohomology!Koszul invariance}
For any same-type finite-type Verdier-Koszul pair
$(\mathcal{A}, \mathcal{A}^!)$ whose module categories are related by
the finite-type Koszul derived equivalence and for which the
Beilinson--Drinfeld chiral Hochschild comparison is Morita-compatible,
there is
an isomorphism of graded rings:
\begin{equation}\label{eq:ch-koszul-iso}
\ChirHoch^*(\mathcal{A}) \;\cong\; \ChirHoch^*(\mathcal{A}^!)
\end{equation}
preserving the curve-level Hilbert profile.  On the critical
formal-disk model of Theorem~\textup{\ref{thm:affine-periodicity-critical}},
the corresponding local statement preserves the number and degrees of
polynomial generators.
In particular, the number of independent deformation directions
$\dim \ChirHoch^2(\mathcal{A})
 = \dim \ChirHoch^2(\mathcal{A}^!)$
is a Koszul-duality invariant.
\end{corollary}

\begin{proof}
The isomorphism~\eqref{eq:ch-koszul-iso} is induced by the module-level
Koszul duality: the derived equivalence
$D^b(\cA\text{-mod}) \simeq D^b(\cA^!\text{-mod})$
(Theorem~\ref{thm:koszul-equivalence-categories}) induces, by
Keller's theorem~\cite{Keller01}, an isomorphism on Hochschild
cohomology for derived-equivalent finite-type categories.  The
Morita-compatible Beilinson--Drinfeld comparison transports this
ordinary Hochschild isomorphism to the chiral Hochschild surface.
This uses the module-level derived equivalence, not bar-cobar
inversion $\Omega(\barB(\cA)) \simeq \cA$; the latter recovers~$\cA$
itself and does not relate~$\cA$ to~$\cA^!$.
Proposition~\ref{prop:periodicity-same-type} gives the same Hilbert
profile on the standard same-type rows.  In the local critical model,
the Gel'fand--Fuchs computation supplies the polynomial generator
statement separately.
\end{proof}

\begin{remark}[Mixed-type pairs]\label{rem:periodicity-mixed-type}
The preceding proposition does not address mixed-type Koszul pairs, where $\mathcal{A}$ and $\mathcal{A}^!$ have structurally different chiral Hochschild cohomology. The computed same-type pairs listed here are $(\widehat{\mathfrak{g}}_k, \widehat{\mathfrak{g}}_{k'})$ and $(\mathrm{Vir}_c, \mathrm{Vir}_{26-c})$; for principal $\mathcal{W}_N$ the same-family statement is proved at the characteristic level and the algebra identification with $\mathcal{W}_N^{k'}$ is conditional on DS/bar transport. The question of whether Koszul duality preserves the chiral Hochschild growth type for genuinely distinct pairs (e.g., a hypothetical Koszul pair between a Kac--Moody algebra and a $\mathcal{W}$-algebra) remains open.
\end{remark}

\begin{remark}[Periodicity profile and upper bounds]\label{rem:periodicity-profile}
\index{periodicity!profile}
For the rank-$1$ local continuous models above, the periodicity
operator is strict ($N \in \{2,4\}$). For rank~$r > 1$, the local
Poincar\'e series
$\prod_{i=1}^r (1 + t^{2m_i+1})/(1 - t^{2(m_i+1)})$
gives polynomial growth
$\dim\,\ChirHoch_{\mathrm{loc}}^n = O(n^{r-1})$, and the lcm
\[
\mathrm{lcm}\bigl(2(m_1{+}1), \ldots, 2(m_r{+}1)\bigr)
\]
is only a quasi-period \emph{upper bound} (no strict period exists for $\mathfrak{sl}_3$; see Remark~\ref{rem:affine-periodicity-analysis}).

\emph{Convention.} $\mathrm{Period}(\mathcal{A})$ is reserved for strict periodicity. The full invariant is the periodicity profile $\Pi(\mathcal{A}) = (M_{\mathcal{A}},\, Q_{\mathcal{A}},\, G_{\mathcal{A}})$ of Remark~\ref{rem:periodicity-triple}.
\end{remark}

\begin{remark}[The periodicity profile]\label{rem:periodicity-triple}
\index{periodicity profile!$\Pi(\mathcal{A})$|textbf}
Define the \emph{periodicity profile}
\[
\Pi(\mathcal{A}) \;=\; \bigl(M_{\mathcal{A}},\; Q_{\mathcal{A}},\;
G_{\mathcal{A}}\bigr)
\]
where:
\begin{enumerate}[label=\textup{(\roman*)}]
\item $M_{\mathcal{A}} = \bigoplus_{n \geq 0} \ChirHoch^n(\mathcal{A})$ is the \emph{periodicity module}, viewed over the periodicity operator $\sigma$. For rank~$1$, $M_{\mathcal{A}}$ is free of rank~$1$ over $\mathbb{Z}[\sigma]/(\sigma^N - 1)$; for rank~$r > 1$, it has $r$~polynomial and $r$~exterior factors.
\item $Q_{\mathcal{A}} = \Lambda(x_{2m_1+1}, \ldots, x_{2m_r+1}) \otimes \mathbb{C}[y_{2(m_1+1)}, \ldots, y_{2(m_r+1)}]$ is the \emph{periodicity algebra} (the Gel'fand--Fuchs ring governing $M_{\mathcal{A}}$).
\item $G_{\mathcal{A}} = \mathrm{Gal}(K_{\mathcal{A}}/\mathbb{Q})$ is the \emph{periodicity Galois group} of the splitting field of the Poincar\'e series denominator. For $\widehat{\mathfrak{sl}}_2$: $G = \mathbb{Z}/2$; for $\widehat{\mathfrak{sl}}_3$: $G$ contains cyclotomic Galois groups of $6$th and $8$th roots.
\end{enumerate}
On the finite-type same-surface rows where Koszul duality acts on
this profile, $M_{\mathcal{A}^!}$ is the graded dual of
$M_{\mathcal{A}}$ (Proposition~\ref{prop:periodicity-exchange-koszul}),
$Q_{\mathcal{A}^!} \cong Q_{\mathcal{A}}$, and
$G_{\mathcal{A}^!} = G_{\mathcal{A}}$.

\textbf{Nilpotence-periodicity interpretation.}
Conjecturally, the profile~$\Pi(\cA)$ is the exponential of the
nilpotent data in $\kappa(\cA)$
(Remark~\ref{rem:nilpotence-periodicity}). The period~$N$ is the
smallest integer with $\exp(2\pi i N \kappa/24) = 1$
(Conjecture~\ref{conj:modular-periodicity-minimal}), and $G_{\cA}$
is the arithmetic shadow of the monodromy representation on bar
cohomology.
\end{remark}

\begin{definition}[Derived scalar period]
\label{def:derived-scalar-period}
\index{periodicity profile!derived scalar period}
\index{scalar period!derived from profile}
A scalar period $\Perbr(\cA) \in \mathbb{Z}_{>0}$ exists if and
only if $M_{\cA}$, $Q_{\cA}$, and $G_{\cA}$ all induce commuting
finite-order automorphisms of a common finite-type complex.
In that case,
\[
\Perbr(\cA)
\;:=\;
\operatorname{ord}\langle T_{\cA},\, R_{\cA},\, G_{\cA} \rangle,
\]
the exponent of the group generated by the modular monodromy
operator~$T_{\cA}$, the spectral $R$-matrix datum~$R_{\cA}$,
and the geometric amplitude operator~$G_{\cA}$.
Otherwise the primary invariant is the full profile
$\Pi_{\cA}$, not an integer.

\smallskip\noindent
\emph{Status.}
For $\widehat{\mathfrak{sl}}_2$ on the strict rank-$1$ periodic
surface, $\Perbr(\cA)$ agrees with the scalar period appearing in the
explicit calculations of Chapter~\ref{chap:deformation-theory}, and
Proposition~\ref{prop:periodicity-exchange-koszul} records its
transport under duality.  For Virasoro on the curve-level chiral
Hochschild surface, Theorem~\ref{thm:virasoro-periodicity} supplies
an amplitude-two class rather than an invertible all-degree period;
an integer period is available only after an external finite-type
completion produces a finite-order automorphism. For
higher-rank algebras where strict periodicity fails
(Remark~\ref{rem:affine-periodicity-analysis}), the derived
period is undefined and $\Pi_{\cA}$ is the correct invariant.
\end{definition}

\begin{proposition}[Numerical admissible data under the level reflection; \ClaimStatusProvedHere]
\label{prop:admissible-levels-permuted}
\index{admissible level!Koszul duality}
Let $\mathfrak{g}$ be a simple Lie algebra. A level $k$ is \emph{admissible}
\textup{(}in the sense of Kac--Wakimoto~\cite{KW88}\textup{)} if it has the form
$k + h^\vee = p/q$ with $p,q$ coprime positive integers,
$p \geq h^\vee$ if $\mathfrak{g}$ is simply laced \textup{(}$p \geq h$ otherwise\textup{)}.
Under the level reflection $k \mapsto k' = -k - 2h^\vee$:
\begin{enumerate}[label=\textup{(\roman*)}]
\item The dual level satisfies $k' + h^\vee = -(k + h^\vee) = -p/q$,
so $|k' + h^\vee| = p/q$ and the denominator is preserved.
\item The Kac--Wakimoto numerator bound on $p$ is preserved as
numerical data.  Literal Kac--Wakimoto admissibility is not preserved
as a statement with $k'+h^\vee>0$; the image is the signed mirror
stratum $k'+h^\vee=-p/q$.
\end{enumerate}
\end{proposition}

\begin{proof}
The level reflection $k \mapsto -k-2h^\vee$ sends
$k + h^\vee \mapsto -(k+h^\vee)$, preserving the absolute value and
negating the sign.  The coprime pair $(p,q)$ and the relevant
Kac--Wakimoto lower bound on~$p$ are unchanged.  Since the definition
of an admissible level uses the positive rational number $p/q$, the
reflected level is a signed mirror point, not a literal positive
admissible level.  Any categorical comparison across the mirror
requires additional representation-theoretic input and is not part of
this proposition.
\end{proof}

\begin{remark}[Module categories are not identified here]
\label{rem:admissible-level-module-frontier}
The proposition above is a numerical statement about the level
involution. A general derived equivalence of admissible module
categories is \emph{not} proved on this surface. The proved
module-level Koszul duality used here is the later
$\Eone$ complete/conilpotent theorem
Theorem~\ref{thm:e1-module-koszul-duality}; extending such a
statement to full admissible chiral module categories remains
frontier material.
\end{remark}

\begin{remark}[The critical locus in level space]\label{rem:critical-locus}
\index{critical level!locus in level space}
The level space $k \in \bC$ is stratified by Koszul behavior:
\begin{center}
\begin{tabular}{ll}
\textbf{Locus} & \textbf{Koszul behavior} \\
\hline
Generic $k$ & Curved affine bar package with level-reflected scalar partner \\
Critical $k = -h^\vee$ & Level fixed point; scalar curvature zero \\
Admissible $k + h^\vee = p/q$ & Positive Kac--Wakimoto stratum; signed mirror under reflection \\
$k \in \bZ_{\geq 0}$ (integrable) & KL equivalence with $U_q(\mathfrak{g})$; $q \leftrightarrow q^{-1}$ \\
$k + h^\vee \in \bZ_{<0}$ & Dual integrable: $k' = -k-2h^\vee \in \bZ_{\geq 0}$
\end{tabular}
\end{center}
The Feigin--Frenkel involution is the reflection
$k + h^\vee \xmapsto{\mathrm{FF}} -(k + h^\vee)$; it preserves the
generic mirror data and has unique fixed point $k=-h^\vee$, while the
positive admissible stratum maps to its signed mirror.
\end{remark}

% ================================================================
% SECTION: FROM QUADRATIC TO NON-QUADRATIC
% ================================================================

\section{The transition from quadratic to non-quadratic Koszul duality}

\subsection{Limitations of quadratic theory}

The Heisenberg algebra is quadratic: a single generator with a single OPE relation (\S\ref{sec:frame-bar-deg1}). Most chiral algebras of physical interest are not.

\begin{definition}[Quadratic chiral algebra]
A chiral algebra $\mathcal{A}$ is \emph{quadratic} if it admits a presentation:
\[
\mathcal{A} = \text{Free}^{\text{ch}}(V)/(R)
\]
where $V$ is a locally free $\mathcal{O}_X$-module concentrated in conformal weight 1, and $R \subset j_*j^*(V \boxtimes V)$ consists of relations among products of two generators.
\end{definition}

\begin{example}[The \texorpdfstring{$\beta\gamma$}{beta-gamma} system is quadratic]
Generators: $\beta$ (weight 1), $\gamma$ (weight 0) \\
Relation: $[\beta(z), \gamma(w)] = \delta(z-w)$ \\
This is quadratic after shifting $\gamma$ to weight 1.
\end{example}

\begin{example}[The Virasoro algebra is non-quadratic]
The stress tensor $T(z)$ has weight 2, and the OPE:
\[
T(z)T(w) = \frac{c/2}{(z-w)^4} + \frac{2T(w)}{(z-w)^2} + \frac{\partial T(w)}{z-w}
\]
Cannot be expressed with only quadratic relations due to the quartic pole.
\end{example}

\subsection{The Maurer--Cartan correspondence for quadratic algebras}

Following Gui, Li, and Zeng~\cite{GLZ22}, we prove a Maurer--Cartan correspondence for quadratic chiral algebras:

\begin{theorem}[Maurer--Cartan correspondence, quadratic case; \ClaimStatusProvedHere]
\label{thm:mc-quadratic}
For a dualizable quadratic chiral algebra $\mathcal{A} = \mathcal{A}(N, P)$ with dual $\mathcal{A}^! = \mathcal{A}(s^{-1}N^{\vee}\omega^{-1}, P^{\perp})$, there is a bijection:
\[
\text{Hom}_{\text{ChirAlg}}(\mathcal{A}, B) \cong \text{MC}(\mathcal{A}^! \otimes B)
\]
where $\text{MC}$ denotes the set of Maurer--Cartan elements.
\end{theorem}

\begin{proof}
\emph{$(\Rightarrow)$}\enspace
Given $\phi\colon \mathcal{A} \to B$, set
$\alpha = (\phi \otimes \mathrm{id})(s^{-1}\mathrm{Id}_{N \otimes N^\vee}) \in (\mathcal{A}^! \otimes B)^1$.
Then $d\alpha = 0$ by $\phi(P) = 0$ and $P^\perp$ orthogonality, and
$[\alpha,\alpha] = 0$ by associativity of $\phi$.

\emph{$(\Leftarrow)$}\enspace
Given $\alpha = \sum_i a_i^! \otimes b_i \in \mathrm{MC}(\mathcal{A}^! \otimes B)$,
define $\phi(n) = \sum_i \langle n, a_i^!\rangle b_i$.
The component $d\alpha = 0$ forces $\phi(P) = 0$ (so $\phi$ factors through $\mathcal{A}$),
and $[\alpha,\alpha] = 0$ forces associativity.
The two constructions are inverse by the universal property of $\mathrm{Free}^{\mathrm{ch}}(N)$. \qedhere
\end{proof}

\subsection{Extending to non-quadratic: higher Maurer--Cartan equations}

For non-quadratic algebras, the quadratic Maurer--Cartan equation is
replaced by the higher equation:

\begin{definition}[\texorpdfstring{$A_\infty$}{A-infinity} Maurer--Cartan equation]
Let $\mathcal{A}$ be a chiral algebra with
$A_\infty$ structure $(m_1, m_2, m_3, \ldots)$.
An element $\alpha \in \mathcal{A}^1$ satisfies the
\emph{$A_\infty$ Maurer--Cartan equation} if:
\[
\sum_{n=1}^{\infty} m_n(\alpha, \ldots, \alpha) = 0
\]
(There are no factorial prefactors: the $m_n$ take ordered inputs, unlike the
symmetric $L_\infty$ brackets $\ell_n$ where the Maurer--Cartan equation
$\sum \frac{1}{n!}\ell_n(\alpha,\ldots,\alpha) = 0$ includes $\frac{1}{n!}$
to account for permutation symmetry.)
\end{definition}

\begin{example}[Cubic relations require \texorpdfstring{$m_3$}{m3}]
For the Yangian with RTT relations (cubic), the MC equation becomes:
\[
m_1(\alpha) + m_2(\alpha, \alpha) + m_3(\alpha, \alpha, \alpha) = 0
\]
where $m_3$ encodes the RTT relation. (In the $L_\infty$ formulation via
antisymmetrization, the corresponding equation is
$d\alpha + \frac{1}{2}[\alpha,\alpha] + \frac{1}{6}\ell_3(\alpha,\alpha,\alpha) = 0$.)
\end{example}

% ================================================================
% SECTION: THE YANGIAN AS NON-QUADRATIC EXAMPLE
% ================================================================

\section{The Yangian: first non-quadratic example}

\subsection{Yangian algebras}

Drinfeld~\cite{Drinfeld85} introduced the Yangian as a canonical deformation of $U(\mathfrak{g}[t])$, motivated by the quantum inverse scattering method and the quantum Yang--Baxter equation.

\subsection{Definition of the Yangian}

\begin{definition}[The Yangian \texorpdfstring{$Y(\mathfrak{g})$}{Y(g)}]\label{def:yangian-kps}
For a simple Lie algebra $\mathfrak{g}$ with basis $\{t_a\}_{a=1}^{\dim \mathfrak{g}}$ and structure constants $[t_a, t_b] = f_{ab}^c t_c$, the Yangian $Y(\mathfrak{g})$ is generated by elements $\{J_n^a : n \geq 0, a = 1, \ldots, \dim \mathfrak{g}\}$ with relations:

\emph{Level-0}: The $J_0^a$ generate a copy of $\mathfrak{g}$:
\[
[J_0^a, J_0^b] = f_{ab}^c J_0^c
\]

\emph{Serre relations}:
\[
[J_0^a, J_n^b] = f_{ab}^c J_n^c
\]

\emph{RTT relations} (the non-quadratic part):
\[
[J_1^a, [J_1^b, J_0^c]] - [J_0^a, [J_1^b, J_1^c]] = \hbar^2 \sum_{d,e,f} A^{abc}_{def}\, \{J_0^d, J_0^e, J_0^f\}
\]
where $A^{abc}_{def} = \tfrac{1}{24}\sum_{p,q,r} f^a_{dp}\, f^b_{eq}\, f^c_{fr}\, f^{pqr}$ and
$\{x,y,z\} = \sum_{\sigma \in S_3} x_{\sigma(1)} y_{\sigma(2)} z_{\sigma(3)}$ is the
symmetrized product (the ``Drinfeld terrible relation''). Higher relations between
$J_r^a$ and $J_s^b$ for $r,s \geq 1$ involve iterated commutators of this basic relation.
\end{definition}

The RTT relations involve products of three generators, making the Yangian inherently non-quadratic.

\subsection{Affine current input and Yangian deformation data}

\begin{theorem}[Affine Kac--Moody as chiral algebra; \ClaimStatusProvedHere]
\label{thm:chiral-yangian-km}
The \emph{affine Kac--Moody algebra} $\hat{\mathfrak{g}}_k$ defines a chiral algebra on $\mathbb{P}^1$, with:
\begin{enumerate}
\item Generating fields $J^a(z) = \sum_{n \in \mathbb{Z}} J_n^a z^{-n-1}$
\item OPE structure:
\[
J^a(z)J^b(w) \sim \frac{f_{ab}^c J^c(w)}{z-w} + \frac{k\, \Omega^{ab}}{(z-w)^2}
\]
where $\Omega^{ab}$ is the Killing form and $k$ is the level
\end{enumerate}
\end{theorem}

\begin{proof}
The chiral algebra axioms for $\hat{\mathfrak{g}}_\kappa$ follow from
the current OPE.

\emph{Locality}: The OPE has poles of order at most~2, ensuring $(z-w)^2 J^a(z)J^b(w) = (z-w)^2 J^b(w)J^a(z)$.

\emph{Associativity}: The Jacobi identity for triple OPEs follows from the Lie algebra Jacobi identity for $\mathfrak{g}$ and the cocycle condition for the central extension.

\emph{Translation covariance}: The $\mathcal{D}_X$-module structure provides the $\partial_z$-action on all fields via $\partial_z J^a(z) = [L_{-1}, J^a(z)]$, where $L_{-1}$ is the infinitesimal translation on~$X$ (part of the intrinsic $\mathcal{D}_X$-module structure, independent of the Sugawara construction). \qedhere
\end{proof}

\begin{theorem}[Critical centre and Yangian deformation data {\cite{Drinfeld85,Feigin-Frenkel}}; \ClaimStatusProvedElsewhere]
\label{thm:chiral-yangian}
The Yangian $Y(\mathfrak{g})$ is Drinfeld's deformation of
$U(\mathfrak{g}[t])$ \cite{Drinfeld85}.  Separately, at the critical
level $\kappa=-h^\vee$, the completed affine Kac--Moody chiral
algebra has Feigin--Frenkel centre
\[
Z(\hat{\mathfrak{g}}_{-h^\vee})
 \cong \mathrm{Fun}(\mathrm{Op}_{\mathfrak{g}^\vee}).
\]
These are distinct external inputs.  This theorem does not assert a
Kazhdan--Lusztig equivalence at critical level between smooth
critical affine modules and Yangian modules; any such comparison
requires a separate quantum-affine/Drinfeld--Kazhdan surface.
\end{theorem}

\begin{remark}[Yangian taxonomy on this surface]
\label{rem:yangian-factorization-scope}
Definition~\ref{def:yangian-kps} is the Drinfeld Yangian
$Y_\hbar(\mathfrak g)$ as an algebra.  A chiral or factorisation
Yangian $Y_X(\mathfrak g)$ is extra structure: it must specify the
$\hbar$-adic topology, the factorisation products on
configuration spaces, and the higher operations whose first
nonquadratic term is the RTT/Serre operation.  The affine current OPE
of Theorem~\ref{thm:chiral-yangian-km} supplies only the binary
current part.  It does not produce the Yangian Hopf coproduct, a
quantum-affine equivalence, or a bialgebra structure on a Koszul dual.
The deconcatenation coproduct on $\bar B(Y_X(\mathfrak g))$ is the
bar-coalgebra coproduct; the Yangian coproduct is a separate
co-operation coming from the tensor product on the chosen
representation category.
\end{remark}

\subsection{Bar complex of the Yangian}

\begin{theorem}[Bar complex structure for a factorized Yangian model; \ClaimStatusConditional]
\label{thm:yangian-bar-complex-structure}
Assume a factorization enhancement $Y_X(\mathfrak g)$ of the
Yangian has been fixed on $X=\mathbb P^1$, with binary current
part given by the affine OPE and first nonquadratic relation given
by Drinfeld's cubic Serre/RTT relation of
Definition~\textup{\ref{def:yangian-kps}}.  Assume also that these
operations are continuous for the $\hbar$-adic/PBW topology and
satisfy the completed $A_\infty$ relations.  Then its ordered chiral
bar complex has degree-$n$ term
\[
\bar{B}^n(Y_X(\mathfrak{g})) =
\Gamma\left(\overline{C}_n(\mathbb{P}^1),
Y_X(\mathfrak{g})^{\boxtimes n} \otimes \Omega^n_{\log}\right),
\]
with differential
$d_{\mathrm{bar}}=d_{\mathrm{int}}+d_2+d_3+\cdots$ whose binary
piece $d_2$ encodes the current Lie bracket and central term, and
whose first higher piece $d_3$ encodes the cubic Serre/RTT relation.
\end{theorem}

\begin{proof}
\emph{Degree 1}: Elements are $J^a(z) \otimes dz$.

\emph{Degree 2}: Elements are $J^a(z_1) \otimes J^b(z_2) \otimes d\log(z_1 - z_2)$.

The differential:
\begin{align}
d(J^a \otimes J^b \otimes \eta_{12}) &= \text{Res}_{z_1 \to z_2} J^a(z_1)J^b(z_2) \otimes \eta_{12} \\
&= f_{ab}^c J^c + \kappa\, \Omega^{ab} \cdot 1
\end{align}

\emph{Degree 3}: Elements $J^a \otimes J^b \otimes J^c \otimes \eta_{12} \wedge \eta_{23}$.

The first higher differential is the image of the cubic relation in
the chosen factorization presentation:
\[
d(\omega_3) = d_2(\omega_3) + d_3(\omega_3),
\qquad
d_3(\omega_3)=\mathrm{RTT}(J^a,J^b,J^c).
\]

Thus nonquadraticity is visible as a higher bar differential on the
completed factorisation model.  The identity
$d_{\mathrm{bar}}^2=0$ is exactly the completed $A_\infty$ relation
among $d_2,d_3,\ldots$; it is an assumption on the chosen
factorisation Yangian model, not a consequence of the affine OPE
alone. \qedhere
\end{proof}



% ================================================================
% SECTION: W-ALGEBRAS
% ================================================================

\section{$\mathcal{W}$-algebras}
\label{sec:w-algebras-kps}

\subsection{Historical development}

The $W_3$ algebra (Zamolodchikov, 1985) arises from quantum Drinfeld--Sokolov reduction (Feigin--Frenkel, 1990).

\subsection{The BRST construction}

\begin{definition}[\texorpdfstring{$\mathcal{W}$}{W}-algebra via quantum Drinfeld--Sokolov]
For a simple Lie algebra $\mathfrak{g}$ and principal nilpotent element $e \in \mathfrak{g}$, the W-algebra $\mathcal{W}^k(\mathfrak{g})$ at level $k$ is:
\[
\mathcal{W}^k(\mathfrak{g}) = H^0_{Q_{DS}}(\hat{\mathfrak{g}}_k \otimes \mathcal{F}_{gh})
\]
where $Q_{DS}$ is the BRST charge and $\mathcal{F}_{gh}$ is the ghost system.
\end{definition}

\begin{example}[\texorpdfstring{$\mathcal{W}_3$}{W3} algebra for \texorpdfstring{$\mathfrak{g} = \mathfrak{sl}_3$}{g = sl3}]
\emph{Step 1}: Start with $\hat{\mathfrak{sl}}_3$ at level $k$.

\emph{Step 2}: Choose principal $\mathfrak{sl}_2$ embedding:
\[
e = \begin{pmatrix} 0 & 1 & 0 \\ 0 & 0 & 1 \\ 0 & 0 & 0 \end{pmatrix}, \quad
h = \begin{pmatrix} 2 & 0 & 0 \\ 0 & 0 & 0 \\ 0 & 0 & -2 \end{pmatrix}, \quad
f = \begin{pmatrix} 0 & 0 & 0 \\ 2 & 0 & 0 \\ 0 & 2 & 0 \end{pmatrix}
\]

\emph{Step 3}: Add ghosts $(b_\alpha, c_\alpha)$ for positive roots.

\emph{Step 4}: BRST charge:
\[
Q = \oint \left( \sum_{\alpha \in \Delta_+} c_\alpha J^\alpha + \text{ghost self-coupling terms} + \text{DS constraint} \right) dz
\]

\emph{Step 5}: Cohomology generators: $T$ (weight 2), $W$ (weight 3).

\emph{Step 6}: The resulting OPEs are given in Definition~\ref{def:w3-algebra}.
The sixth-order pole in the $W$--$W$ OPE makes this highly non-quadratic.
\end{example}

\subsection{Critical centre and the screening model}

\begin{theorem}[Feigin--Frenkel centre {\cite{FF}}; \ClaimStatusProvedElsewhere]
\label{thm:feigin-frenkel-bar}
At critical level $k=-h^\vee$, the centre of the completed affine
chiral algebra is
\[
\mathfrak z(\widehat{\mathfrak g})\cong
\mathrm{Fun}(\mathrm{Op}_{{}^L\!\mathfrak g}),
\]
the algebra of functions on Langlands-dual opers.  Under the Miura
free-field realization of the principal $\mathcal W$-algebra, the
screening operators $S_1,\ldots,S_r$ describe the intersection of
screening kernels.
\end{theorem}

\begin{proof}
This is the Feigin--Frenkel theorem at the critical level, together
with the standard Miura/screening realization of principal
$\mathcal W$-algebras.  The statement is about the centre and the
free-field screening presentation; it is not a theorem that the
chiral bar coalgebra of $\mathcal W^{-h^\vee}(\mathfrak g)$ is equal
to a polynomial algebra on screenings.
\end{proof}

\begin{remark}[Screening Koszul complex as model]
\label{rem:w-critical-bar-viewpoint}
The complex
\[
\mathrm{Sym}[S_1,\ldots,S_r]\otimes\Omega^*_{\log}
\]
is a useful associated-graded or free-field model for the
critical-level screening calculation.  It may serve as a lead for a
bar computation after the Miura comparison, logarithmic extension,
and completion hypotheses are supplied.  This chapter does not use it
as an identification of $\bar B(\mathcal W^{-h^\vee}(\mathfrak g))$.
\end{remark}

\subsection{Langlands duality for $\mathcal{W}$-algebras}

\begin{remark}[Frenkel--Gaitsgory viewpoint {\cite{FG06}}; \ClaimStatusHeuristic]
\label{rem:frenkel-gaitsgory-langlands-viewpoint}
At critical level, the local geometric/quantum-Langlands literature
relates affine and oper-theoretic data for $\mathfrak g$ and its
Langlands dual $\mathfrak g^L$. This critical-level picture supplies
external comparison data for principal $\mathcal W$-algebras and
Langlands-dual data. It is
\emph{not} identified with a proved bar-cobar Koszul-duality theorem of
the form
\[
\mathcal{W}^{-h^{\vee}}(\mathfrak{g})^! \simeq
\mathcal{W}^{-h^{\vee}}(\mathfrak{g}^L),
\]
and no such theorem is claimed internally here.
\end{remark}

% ================================================================
% SECTION: NON-PRINCIPAL W-ALGEBRAS
% ================================================================

\section{Non-principal $\mathcal{W}$-algebras}

\subsection{Non-principal orbits}

Quantum Drinfeld--Sokolov reduction attached to a non-principal
nilpotent orbit produces a non-principal $\mathcal W$-algebra.  The
exchange of non-principal orbits under Gaiotto--Witten S-duality is
a conjectural external duality input unless the orbit and level
transform are supplied separately.

\subsection{\texorpdfstring{Example: Subregular $\mathcal{W}$-algebra for $\mathfrak{sl}_4$}{Example: Subregular $\mathcal{W}$-algebra for sl 4}}

\begin{definition}[Subregular nilpotent]
The subregular nilpotent in $\mathfrak{sl}_4$ has Jordan type $(3,1)$:
\[
e_{subreg} = \begin{pmatrix} 0 & 1 & 0 & 0 \\ 0 & 0 & 1 & 0 \\ 0 & 0 & 0 & 0 \\ 0 & 0 & 0 & 0 \end{pmatrix}
\]
\end{definition}

\begin{theorem}[Structure of \texorpdfstring{$\mathcal{W}(\mathfrak{sl}_4, e_{subreg})$}{W(sl4, e\_subreg)} {\cite{KRW}}; \ClaimStatusProvedElsewhere]\label{thm:w-algebra-sl4}
The subregular W-algebra has generators $T$ (weight~$2$), $G^{\pm}$ (weight~$3/2$), and $J$ (weight~$1$), with OPEs involving fractional powers and fermionic statistics.
\end{theorem}

The fractional weights require orbifold constructions on configuration spaces.

\subsection{S-duality and Koszul duality}

\begin{theorem}[Feigin--Frenkel duality as S-duality, principal simply-laced case; \ClaimStatusProvedElsewhere]\label{thm:ff-s-duality}
For simply-laced $\mathfrak{g}$ and principal nilpotent $f = f_{\mathrm{prin}}$,
the Gaiotto--Witten S-duality reduces to Feigin--Frenkel duality
\cite{Feigin-Frenkel}:
\[
\mathcal{W}^k(\mathfrak{g}, f_{\mathrm{prin}})
\;\simeq\;
\mathcal{W}^{k'}(\mathfrak{g}, f_{\mathrm{prin}}),
\qquad (k + h^{\vee})(k' + h^{\vee}) = 1.
\]
The duality inverts the shifted level and preserves the central
charge; for $\mathfrak{g} = \mathfrak{sl}_2$ it is the Virasoro
self-duality $b \mapsto b^{-1}$ with $b^{2} = k+2$, under which
$c = 13 - 6(b^{2} + b^{-2})$ is fixed. It is distinct from the
reflection $k \mapsto -k - 2h^{\vee}$, which sends $k + h^{\vee}$ to
$-(k + h^{\vee})$ and hence $\kappa$ to $-\kappa$
(for Virasoro $c + c' = 26$):
that reflection is the curvature-pairing involution, not an
isomorphism of $\mathcal{W}$-algebras.
\end{theorem}

\begin{conjecture}[Gaiotto-Witten S-duality, general case; \ClaimStatusConjectured]\label{conj:gw-s-duality}
For general $\mathfrak{g}$ and nilpotent $f$, there exists a duality:
\[
\mathcal{W}^k(\mathfrak{g}, f) \longleftrightarrow \mathcal{W}^{k^L}(\mathfrak{g}^L, f^L)
\]
where:
\begin{itemize}
\item $k^L = -h^{\vee}(\mathfrak{g}^L) + r^{\vee}/(k + h^{\vee}(\mathfrak{g}))$, where $r^{\vee}$ is the lacing number of $\mathfrak{g}$
\item $f^L$ is the Spaltenstein dual nilpotent
\end{itemize}
\end{conjecture}

\begin{remark}[Scope]
The principal simply-laced case is Feigin--Frenkel
\cite{Feigin-Frenkel}. In type~$A$, the hook-type transport
corridor is a conditional theorem whose missing hypothesis is
uniform DS/bar compatibility over the hook family
\cite{FehilyHook,CLNS24,CFLN24,GenraStages,ButsonNair,AvE23}
(Remark~\ref{rem:hook-type-evidence},
Theorem~\ref{thm:hook-transport-corridor}). Kac--Roan--Wakimoto
\cite{KRW} supplies DS reduction for arbitrary~$f$; the separate
Koszul open mathematical condition is compatibility of Verdier-Koszul duality
with arbitrary-nilpotent reduction, together with the
orbit-dependent level transform.
\end{remark}

Non-principal DS reduction produces non-quadratic chiral algebras in
every type.  Producing Koszul dual pairs from them requires the
additional DS/bar compatibility and orbit-dependent level transform
named above.

% ================================================================
% SECTION: MODULE CATEGORIES AND APPLICATIONS
% ================================================================

\section{Module categories and resolutions}

\subsection{Derived and exotic categorical comparisons}

\begin{theorem}[Koszul equivalence of categories {\cite{BGS96}}; \ClaimStatusProvedElsewhere]
\label{thm:koszul-equivalence-categories}
Let $(\mathcal{A},\mathcal{A}^i)$ be a quadratic, uncurved,
finite-type Koszul pair on the genus-$0$ $\Eone$ surface, and assume
that the linear dual algebra
$\mathcal{A}^!=(\mathcal{A}^i)^\vee$ exists and that the
bounded-below module categories are generated by finite-length
Koszul modules.  Then there is an equivalence of triangulated
categories:
\[
D^b(\mathcal{A}\text{-mod}) \simeq D^b(\mathcal{A}^!\text{-mod})
\]
\end{theorem}

\begin{remark}\label{rem:module-equiv-caveat}
The theorem is not a definition of chiral Koszulness.  In the general
chiral setting \textup{(}curved, non-quadratic, or only completed
finite type\textup{)}, the proved intrinsic statement is the
coderived/contraderived comparison for the bar coalgebra
$C=\barB^{\mathrm{ch}}(\cA)$ and its completed-dual module category
\textup{(}Theorems~\ref{thm:positselski-chiral} and
\ref{thm:full-derived-module-equiv}\textup{)}. The ordinary derived
module equivalence is recovered only on the flat finite-type bar-dual
surface, including the quadratic genus-$0$ setting of
Theorem~\ref{thm:e1-module-koszul-duality}.
\end{remark}

\begin{theorem}[Positselski equivalence for the finite-type chiral bar
coalgebra; \ClaimStatusConditional]
\label{thm:positselski-chiral}
\index{Positselski!comodule-contramodule}
\index{derived category!comodule-contramodule equivalence}
For a Koszul chiral algebra $\mathcal{A}$ whose curved bar complex
$C=\bar{B}^{\mathrm{ch}}(\mathcal{A})$ lies in the finite-type
conilpotent CDG surface and satisfies the product/sum exactness
hypothesis of
Theorem~\textup{\ref{thm:chiral-co-contra-correspondence}}, there is
an equivalence of
triangulated categories:
\begin{equation}\label{eq:positselski-equivalence}
D^{\mathrm{co}}(C\text{-}\mathrm{comod}^{\mathrm{conil,\,ch}})
\;\simeq\;
D^{\mathrm{ctr}}(C\text{-}\mathrm{contra}^{\mathrm{ch}})
\end{equation}
between the coderived category of conilpotent chiral
CDG-comodules over the bar coalgebra~$C$ and the contraderived
category of chiral CDG-contramodules over~$C$. In finite type the
right-hand side may be rewritten as complete modules over the graded
dual algebra~$C^*$, and on the quadratic genus-$0$ bar-dual surface
this recovers the module Koszul duality of
Theorem~\textup{\ref{thm:e1-module-koszul-duality}}.
For class-$\mathsf M$ raw direct-sum bar coalgebras the applicable
surface is the strict completed one of
Theorem~\textup{\ref{thm:chiral-positselski-weight-completed}}, with
its explicit continuous co/contra hypotheses.
\end{theorem}

\begin{proof}
This is Theorem~\ref{thm:positselski-chiral-proved}, which follows
from the chiral comodule-contramodule correspondence
(Theorem~\ref{thm:chiral-co-contra-correspondence}) applied to the
CDG-coalgebra $C = \bar{B}^{\mathrm{ch}}(\cA)$. The foundational
chiral coalgebra homological algebra is developed in
\S\ref{sec:chiral-coalgebra-homalg}.
\end{proof}

\begin{remark}[Three comparison layers]
\label{rem:positselski-proof-strategy}
\index{Positselski!comparison layers}
Three layers: (1)~$\bar{B}^{(g)}(\cA)$ is a chiral CDG-coalgebra
\textup{(}Definition~\ref{def:chiral-CDG-coalgebra}\textup{)};
curvature at $g \geq 1$ forces the coderived/contraderived framework
\textup{(}Proposition~\ref{prop:curved-bar-acyclicity}\textup{)}.
(2)~The functors $\Phi_C^{\mathrm{ch}}, \Psi_C^{\mathrm{ch}}$
\textup{(}Definition~\ref{def:chiral-Phi-Psi}\textup{)} descend to
$D^{\mathrm{co}} \simeq D^{\mathrm{ctr}}$
\textup{(}Theorem~\ref{thm:chiral-co-contra-correspondence}\textup{)}.
(3)~In finite type, CDG-contramodules over
$\bar{B}^{\mathrm{ch}}(\cA)$ may be rewritten as complete modules over
the graded dual algebra $\bar{B}^{\mathrm{ch}}(\cA)^*$; the further
identification with $\cA^!$ is only available on the quadratic
genus-$0$ bar-dual surface
\textup{(}Theorem~\ref{thm:bar-computes-dual}\textup{)}.
\end{remark}

\begin{theorem}[Flat finite-type reduction on the completed-dual side; \ClaimStatusConditional]
\label{thm:full-derived-module-equiv}
\index{derived category!flat finite-type reduction}
\index{Positselski!derived reduction on completed dual}
Let $C=\barB^{\mathrm{ch}}(\cA)$ be the chiral bar coalgebra of a
Koszul chiral algebra~$\cA$, and assume we are on a flat finite-type
locus where Positselski's comparison theorem identifies the
contraderived category of chiral CDG-contramodules over~$C$ with the
ordinary bounded derived category of complete modules over the graded
dual algebra~$C^*$. Then there is a triangulated equivalence:
\begin{equation}\label{eq:full-derived-equiv}
D^{\mathrm{ctr}}(C\text{-}\mathrm{contra}^{\mathrm{ch}})
\;\simeq\;
D^b(\mathrm{Mod}^{\mathrm{compl}}(C^*))
\end{equation}
Consequently, on that same locus, the intrinsic coderived/contraderived
equivalence of Theorem~\textup{\ref{thm:positselski-chiral}} reduces to
an ordinary derived comparison on the completed-dual side. On the
quadratic genus-$0$ bar-dual surface where $C^*\cong\cA^!$, this
recovers Theorem~\textup{\ref{thm:e1-module-koszul-duality}}.
\end{theorem}

\begin{proof}
This is Theorem~\ref{thm:full-derived-module-equiv-proved}.
It is a flat finite-type reduction theorem for the completed-dual side
of the intrinsic Positselski equivalence. The quadratic genus-$0$
specialization then recovers
Theorem~\ref{thm:e1-module-koszul-duality}. See
\S\ref{subsec:positselski-chiral-equivalence} for full details.
\end{proof}

\subsection{Explicit resolutions for non-quadratic cases}

\begin{example}[BGG resolution for W-algebras]
On the external DS-reduction surface for admissible
$\mathcal{W}^k(\mathfrak{g})$-modules, the BGG-type resolution
problem for a simple module $L(\lambda)$ has the form
\[
\cdots \to \bigoplus_{\ell(w)=2} M(w \cdot \lambda) \to \bigoplus_{\ell(w)=1} M(w \cdot \lambda) \to M(\lambda) \to L(\lambda) \to 0
\]
where $w$ ranges over elements of the Weyl group of length $\ell(w)$,
$w \cdot \lambda = w(\lambda + \rho) - \rho$ is the dot action, and
the differentials are given by screening operators. On the
claim used here, the relevant exactness input is
the external exactness of Drinfeld--Sokolov reduction on the
appropriate module category \textup{(}Arakawa~\cite{Arakawa17}\textup{)},
not rationality of $\mathcal{W}^k(\mathfrak{g})$ by itself.
\end{example}

\begin{remark}[Yangian module resolutions require Drinfeld--Kazhdan input]
\label{rem:yangian-module-resolution-scope}
On an external Drinfeld--Kazhdan comparison surface, a resolution
theorem for finite-dimensional $Y(\mathfrak{g})$-modules would have
to construct complexes induced from the quantum affine side,
schematically of the form
\[
\cdots \to Y \otimes V_2 \to Y \otimes V_1 \to Y \otimes V_0 \to M \to 0
\]
where the $V_i$ come from the quantum affine category and the
differentials are encoded by the relevant $R$-matrix data.  This
requires compact generation by the chosen quantum-affine objects,
exactness of the comparison functor, and compatibility of the
braiding reversal $R_q\mapsto R_{q^{-1}}$ with the bar filtration.
The bar-cobar counit $\Omega\bar B(Y_X)\to Y_X$ supplies none of
these three inputs.
\end{remark}

% ================================================================
% SECTION: DEFORMATION THEORY
% ================================================================

\section{Deformation theory and Maurer--Cartan elements}

\subsection{Deforming chiral algebras}

\begin{remark}[Deformation-theoretic scope]
\label{rem:nonquadratic-deformation-scope}
For a non-quadratic chiral algebra, the deformation-theoretic data on
the surface used here are organized by the antisymmetrized
bar $L_\infty$ package discussed below. Maurer--Cartan equations
there record candidate formal deformations.  No stand-alone
comparison theorem is proved here identifying formal deformations of
an arbitrary chiral algebra with
Maurer--Cartan elements of a specifically defined chiral Hochschild
complex; nor is a proved $\beta\gamma$ deformation family
constructed from an explicit Maurer--Cartan element.
\end{remark}

% ================================================================
% SECTION: THE CHERN-SIMONS CONNECTION
% ================================================================

\section{The Chern--Simons structure in non-quadratic Koszul duality}

\subsection{Maurer--Cartan and Chern--Simons}

The Maurer--Cartan equation for non-quadratic chiral algebras:
\[
d\alpha + \frac{1}{2}\ell_2(\alpha, \alpha) + \frac{1}{6}\ell_3(\alpha, \alpha, \alpha) + \cdots = 0
\]
(using the $L_\infty$ brackets $\ell_n$ obtained by antisymmetrizing the $A_\infty$ operations)
is the basic algebraic deformation equation attached to the
antisymmetrized bar $L_\infty$ structure. In quadratic cases it
reduces to the usual dg~Lie Maurer--Cartan equation; any comparison
with Chern--Simons theory requires additional external input.

\subsection{Witten's Chern--Simons interpretation}

Chern--Simons theory on a 3-manifold $M$ with gauge group $G$ at level $k$ produces~\cite{Wit89} the WZW model on $\partial M$ and quantum group representations at $q = e^{\pi i/(k+h^\vee)}$. The action is:
\[
S_{CS}(A) = \frac{k}{4\pi} \int_M \text{Tr}\left(A \wedge dA + \frac{2}{3}A \wedge A \wedge A\right)
\]

\subsection{The precise connection}

\begin{theorem}[Affine Kac--Moody Maurer--Cartan and curvature package;
\ClaimStatusConditional]\label{thm:cs-koszul-km}
\index{Chern--Simons!from Koszul duality}
For a simple Lie algebra $\fg$ and its affine extension
$\widehat{\fg}_k$ at $k \neq -h^\vee$:
\begin{enumerate}[label=\textup{(\roman*)}]
\item \emph{Classical Maurer--Cartan = flatness.}
 The ordinary dg~Lie algebra $\fg \otimes \Omega^*(M)$ carries the
 classical Maurer--Cartan equation
 \[
 d\alpha + \tfrac{1}{2}[\alpha, \alpha] = 0
 \]
 for a $3$-manifold $M$, and this is the flatness equation
 $F_A = 0$. Separately, the Chevalley--Eilenberg complex
 $C^*(\fg)$ is the bar construction $\barB(U\fg)$.
\item \emph{Affine bar = semi-infinite complex.}
 The bar--semi-infinite quasi-isomorphism
 \[
 \barB(\widehat{\fg}_k) \simeq
 C^*_{\infty/2}(\widehat{\fg}_k)
 \]
 \textup{(}Theorem~\ref{thm:bar-semi-infinite-km}\textup{)}
 identifies the affine bar complex with the semi-infinite
 complex of the vacuum module.
\item \emph{Genus-$1$ curvature coefficient and dual cancellation.}
 At genus~$1$, the bar complex has $\dfib^{\,2} = \kappa \cdot \omega_1$
 with $\kappa = (k+h^\vee)\dim(\fg)/(2h^\vee)$. This curvature is
 algebraically nonzero away from the critical level, and for the
 affine Kac--Moody dual pair $(k, -k-2h^\vee)$ one has
 $\kappa(\widehat{\fg}_k) + \kappa(\widehat{\fg}_{k'}) = 0$
 \textup{(}this anti-symmetry is specific to affine Kac--Moody;
 for Virasoro, $\kappa + \kappa' = 13$\textup{)}.
\end{enumerate}
\end{theorem}

\begin{proof}
Part~(i) is classical: the MC equation
$d\alpha + \frac{1}{2}[\alpha, \alpha] = 0$ in the dg~Lie algebra
$\fg \otimes \Omega^*(M)$ is the flatness equation $F_A = 0$ by
definition. Independently, $C^*(\fg) = \barB(U\fg)$ by the
Chevalley--Eilenberg theorem and provides the standard cochain
model attached to $\fg$. The flatness equation is the
Euler--Lagrange equation for $S_{\mathrm{CS}}$: varying
$\delta S_{\mathrm{CS}}/\delta A = \frac{k}{2\pi} F_A = 0$.

Part~(ii) is exactly Theorem~\ref{thm:bar-semi-infinite-km}.

Part~(iii) follows from
$\dfib^{\,2} = \kappa \cdot \omega_1$
(Theorem~\ref{thm:genus-universality}) with
$\kappa(\widehat{\fg}_k) = (k+h^\vee)\dim(\fg)/(2h^\vee)$
(Remark~\ref{rem:km-central-charge-sum}). For the dual pair
$(k, k' = -k-2h^\vee)$,
$\kappa(\widehat{\fg}_k) + \kappa(\widehat{\fg}_{k'}) = 0$
\textup{(}Corollary~\ref{cor:anomaly-duality-km-vol1};
the vanishing is an affine Kac--Moody identity\textup{)}.
\end{proof}

\begin{remark}[External physical interpretation]
\label{rem:cs-km-physical}
Theorem~\ref{thm:cs-koszul-km} supplies algebraic input often
compared with the Chern--Simons/WZW picture: classical
Maurer--Cartan flatness in $\fg \otimes \Omega^*(M)$, the affine bar/semi-infinite
identification, and the dual curvature coefficient. Turning this
into a full boundary Chern--Simons interpretation uses additional
external results and is not proved here.
\end{remark}

\begin{theorem}[\texorpdfstring{$L_\infty$}{L-infinity} Maurer--Cartan equation from a transferred \texorpdfstring{$A_\infty$}{A-infinity} model; \ClaimStatusProvedHere]
\label{thm:linf-mc-flatness}
\index{Chern--Simons!Koszul correspondence}
\index{L-infinity algebra@$L_\infty$ algebra!Maurer--Cartan}
For a chiral algebra $\cA$ with $A_\infty$ operations
$\{m_n\}$ on an underlying graded space, the antisymmetrized
operations $\ell_n$ define an $L_\infty$ structure on that same
graded space. For W-algebras this applies only after a transferred
model has been supplied. The Maurer--Cartan
equation
\[
d\alpha + \tfrac{1}{2}\ell_2(\alpha, \alpha)
+ \tfrac{1}{6}\ell_3(\alpha, \alpha, \alpha)
+ \cdots = 0
\]
is the Maurer--Cartan equation of this transferred
$L_\infty$ algebra.
\begin{enumerate}[label=\textup{(\roman*)}]
\item For quadratic $\cA$
 \textup{(}Kac--Moody\textup{)}: $\ell_n = 0$ for
 $n \geq 3$, and the MC equation is the standard
 dg~Lie Maurer--Cartan equation.
\item For general non-quadratic $\cA$: the higher
 $\ell_n$ are determined by the higher bar operations
 $m_n$ and contribute the corresponding higher terms to the
 Maurer--Cartan equation.
\end{enumerate}
\end{theorem}

\begin{proof}
The $L_\infty$ structure on the underlying graded space of the
$A_\infty$ model is obtained by the standard
antisymmetrization (Loday--Vallette \cite[Thm.~10.1.13]{LV12}):
$\ell_n(x_1, \ldots, x_n) = \sum_{\sigma \in S_n}
\epsilon(\sigma)\, m_n(x_{\sigma(1)}, \ldots,
x_{\sigma(n)})$. The MC equation
$\sum_{n \geq 1} \frac{1}{n!}\ell_n(\alpha^{\otimes n})
= 0$ is then automatic from the $A_\infty$
relations (which are equivalent to the $L_\infty$
Jacobi identities under antisymmetrization).

For~(i): when $\cA$ is quadratic (Kac--Moody), the
transferred model has $m_n = 0$ for $n \geq 3$
(the OPE has only simple and double poles), so
$\ell_n = 0$ for $n \geq 3$ and the MC equation
reduces to $d\alpha + \frac{1}{2}[\alpha, \alpha] = 0$,
which is the usual dg~Lie Maurer--Cartan equation.

For~(ii): if some higher bar operations $m_n$ are nonzero, then
their antisymmetrizations contribute nonzero higher brackets
$\ell_n$, and the Maurer--Cartan equation acquires the displayed
higher terms.
\end{proof}

\label{thm:cs-koszul-general}%
\label{eq:deformed-cs-general}%

\begin{remark}[Algebraic and physical interpretations]
\label{rem:cs-koszul-external}
Theorem~\ref{thm:linf-mc-flatness} establishes the algebraic statement:
the transferred
$A_\infty$ models determine antisymmetrized $L_\infty$ structures,
and their Maurer--Cartan equations
have the displayed higher terms.  The additional assertion that
these terms are flatness equations, BV-BRST vertices, or
Chern--Simons corrections is external: it requires geometric fields,
a cyclic BV pairing, and perturbative QFT input not constructed here.
\end{remark}

\begin{remark}[Scope]
Theorem~\ref{thm:linf-mc-flatness} constructs the Maurer--Cartan
equation on the transferred $L_\infty$ model associated to a chosen
$A_\infty$ structure; Theorem~\ref{thm:cs-koszul-km}
records the affine Kac--Moody classical/semilinear package
\textup{(}classical Maurer--Cartan flatness, bar/semi-infinite
comparison, and the dual curvature coefficient\textup{)};
the displayed higher Maurer--Cartan equation records the
non-quadratic case. Any
field-theoretic Chern--Simons interpretation remains external
to the proved algebraic core.
\end{remark}

\subsection{Quantum-group and Chern--Simons comparison data}

\begin{conjecture}[Witten--Reshetikhin--Turaev comparison surface;
\ClaimStatusConjectured]\label{conj:wrt-conjecture}
Assume:
\begin{enumerate}[label=\textup{(W\arabic*)}]
\item the Witten--Reshetikhin--Turaev construction identifies
the Chern--Simons boundary category at level~$k$ with the
corresponding quantum-group representation category at
\[
 q=\exp\!\left(\frac{\pi i}{k+h^\vee}\right);
\]
\item a Drinfeld--Kazhdan/Kazhdan--Lusztig comparison is fixed on an
evaluation-generated Yangian/quantum-affine surface;
\item the finite-type Verdier-Koszul comparison is available on that
same surface and sends the level parameter by
$k+h^\vee\mapsto -(k+h^\vee)$.
\end{enumerate}
Then the boundary quantum parameter is sent to
\[
 q \longmapsto q^{-1}.
\]
The resulting datum compares the WRT boundary category, the external
Yangian/quantum-affine comparison, and the finite-type
Verdier-Koszul reflection.  It is not a construction of a
Chern--Simons action from the bar coalgebra and not an identification
of $Y_\hbar(\mathfrak g)$ with the Koszul dual of
$V_k(\mathfrak g)$.
\end{conjecture}

\begin{remark}[Scope]
The transferred Maurer--Cartan equation on the chosen
$L_\infty$ model has the formal shape of a deformed flatness
equation \textup{(}preceding theorem\textup{)}.  A physical
Chern--Simons interpretation requires the WRT package
\cite{Wit89,RT91}, a cyclic BV pairing, and a perturbative
quantization theorem.  The algebraic Koszul package supplies only
the boundary bar, the Verdier partner when finite type is present,
and the $L_\infty$ Maurer--Cartan equation.
\end{remark}

If the transferred $L_\infty$ model is cyclic with pairing
$\langle-,-\rangle_{\mathrm{cyc}}$, the first RTT contribution to a
formal action has the typed form
\[
S_{\mathrm{form}}(\alpha)
=\frac12\langle\alpha,d\alpha\rangle_{\mathrm{cyc}}
+\frac16\langle\alpha,\ell_2(\alpha,\alpha)\rangle_{\mathrm{cyc}}
+\frac1{24}
\langle\alpha,\ell_3^{\mathrm{RTT}}(\alpha,\alpha,\alpha)
\rangle_{\mathrm{cyc}}+\cdots .
\]
Without the cyclic pairing and BV integration data this expression is
only a formal cyclic $L_\infty$ functional.

\subsection{Examples of Chern--Simons structure}

\subsubsection{Yangian deformation datum}

For a factorisation Yangian model whose first higher bracket is
$\ell_3^{\mathrm{RTT}}$, the Maurer--Cartan equation is
\[
d\alpha + \frac{1}{2}[\alpha, \alpha]
+ \frac{1}{6}\ell_3^{\mathrm{RTT}}(\alpha, \alpha, \alpha) = 0.
\]
This is an algebraic deformation of the classical Maurer--Cartan
equation inside the chosen completed $L_\infty$ algebra.  It becomes
a gauge-field equation only after an external Chern--Simons/BV model
realizes the same $L_\infty$ algebra as fields.

\subsubsection{Critical-level \texorpdfstring{$\mathcal{W}$}{W} screening datum}

At critical level $k = -h^\vee$:
\[
d\alpha + \sum_{i=1}^r S_i(\alpha) = 0
\]
where $S_i$ are screening charges in the free-field model.  This
equation belongs to the Miura/screening complex; a field-theoretic
identification requires a separate BRST/BV comparison.

\subsubsection{Non-principal \texorpdfstring{$\mathcal{W}$}{W} orbifold datum}

For fractional gradings, an orbifold Chern--Simons model would first
require a genuine orbifold field theory on $M/\Gamma$.  The notation
\[
S = \frac{1}{|\Gamma|} \int_{M/\Gamma} S_{\mathrm{CS}}(\alpha)
\]
records only the orbifold volume normalization by the group
$\Gamma$; no action functional on the non-principal
$\mathcal W$-algebra is constructed from it here.

\subsection{The holographic interpretation}

\begin{conjecture}[AdS/CFT comparison datum; \ClaimStatusConjectured]
\label{conj:ads-cft-cs-koszul}
Assume:
\begin{enumerate}[label=\textup{(H\arabic*)}]
\item a $3$d Chern--Simons or holomorphic-topological bulk theory on
AdS$_3$ has been constructed;
\item its boundary category is identified with modules for a boundary
chiral algebra \textup{(}WZW, Virasoro, or $\mathcal W$\textup{)};
\item the boundary chiral algebra lies on the finite-type
Verdier-Koszul surface of this chapter;
\item a holographic comparison functor identifies physical bulk
observables with the intrinsic Hochschild closed-sector slot
$Z^{\mathrm{der}}_{\mathrm{ch}}(\cA)$.
\end{enumerate}
Then the boundary side separates into three typed objects:
bar-cobar inversion recovers the boundary algebra~$\cA$, Verdier
duality supplies the Koszul partner~$\cA^!$ when finite type is
present, and chiral Hochschild cochains supply the intrinsic
closed-sector operator complex.  The conjecture is the existence of the
external comparison functor in \textup{(H4)}; bulk reconstruction is
not a theorem of bar-cobar inversion.
\end{conjecture}

\begin{remark}[Scope]
The AdS$_3$/CFT$_2$ interpretation requires: (a)~3d gravity as
$SL(2,\mathbb{R})^2$ CS (Witten \cite{Wit89}); (b)~boundary WZW
identification (Beilinson--Drinfeld \cite{BD04}); (c)~a holographic
dictionary comparing the intrinsic Hochschild closed-sector slot with physical
bulk observables. The proved input is the boundary-side
Koszul/module package; bulk reconstruction is not proved here.
\end{remark}

The external boundary dictionary has the type form
\[
\begin{array}{ccc}
\text{bulk Chern--Simons package on AdS}_3
& \rightsquigarrow_{\mathrm{ext}} &
\text{CFT}_2 \text{ on } \partial(\text{AdS}_3) \\
\text{flat connections} & \rightsquigarrow_{\mathrm{ext}} &
\text{chiral algebra modules} \\
\text{Wilson lines} & \rightsquigarrow_{\mathrm{ext}} &
\text{vertex operators} \\
\text{MC elements} & \rightsquigarrow_{\mathrm{alg}} &
\text{formal deformations}
\end{array}
\]
Only the last algebraic arrow is supplied by the transferred
$L_\infty$ package in this chapter.

\subsection{BV formalism}

The BV algebra structure on the bar complex is the fifth piece of the
master differential
(Theorem~\ref{thm:convolution-dg-lie-structure}).

\begin{theorem}[BV structure on bar complex; \ClaimStatusConditional]
\label{thm:bv-structure-bar}
\index{BV algebra!on bar complex}
\index{BV operator!from non-separating clutching}
The chiral bar complex carries a BV algebra structure where:
\begin{enumerate}
\item The BV operator $\Delta$ is the non-separating clutching map
of the modular operad: for an element at genus~$g$ with $n+2$
inputs, $\Delta$ contracts a pair of inputs using the invariant
form $\langle -, - \rangle_\cA$ and produces an element at
genus~$g+1$ with $n$ inputs. This is the $\hbar\Delta$ term in
the five-component differential~\eqref{eq:D-five-components}.
\item The antibracket $\{-,-\}$ is the separating clutching:
joining two elements by identifying one output with one input
through the propagator $P_\cA$. This is the Lie bracket of the
modular convolution algebra
(Theorem~\ref{thm:convolution-dg-lie-structure}).
\item The quantum master equation
$D\Theta^{\leq r} + \tfrac12[\Theta^{\leq r},\Theta^{\leq r}] = 0$
holds at each finite order of the shadow obstruction tower
\textup{(}Remark~\textup{\ref{rem:unifying-principle})}.
\end{enumerate}

\emph{Analytic realization.} On configuration spaces,
$\Delta$ acts by residue extraction along collision diagonals and
$\{-,-\}$ by the symplectic pairing on $T^*C_n(X)$. The Arnold
relation
$\eta_{ij} \wedge \eta_{jk} + \eta_{jk} \wedge \eta_{ki}
+ \eta_{ki} \wedge \eta_{ij} = 0$ is the genus-$0$ shadow of
$\Delta^2 = 0$.
\end{theorem}

\begin{proof}
The BV operator $\Delta$ is the non-separating gluing map
$\xi_{\mathrm{ns}}\colon
\overline{\mathcal{M}}_{g,n+2} \to \overline{\mathcal{M}}_{g+1,n}$
of the modular operad
(Definition~\ref{def:modular-operad}), composed with the invariant
pairing. The antibracket is the separating gluing. The second-order
property of $\Delta$ (the BV axiom) follows from the
codimension-$2$ boundary relations on
$\overline{\mathcal{M}}_{g,n}$: two non-separating degenerations
compose to give a codimension-$2$ stratum that also arises from
a separating degeneration, yielding the BV relation
$\Delta^2 = 0$ and $\Delta(ab) = \Delta(a)b \pm a\Delta(b)
\pm \{a,b\}$. The QME is the Maurer--Cartan equation in the
modular convolution algebra. At the full homotopy level, the BV
structure is the dg shadow of the quantum $L_\infty$-algebra of
Theorem~\ref{thm:modular-quantum-linfty}: $\Delta$ is $\ell_1^{(1)}$
and $\{-,-\}$ is $\ell_2^{(0)}$.
Analytically, this reproduces the
residue calculus on $\overline{C}_n(X)$ by the Prism Principle
(Theorem~\ref{thm:prism-operadic}).
\end{proof}

\begin{remark}[Chern--Simons interpretation]\label{rem:cs-interp}
The BV antibracket on the bar complex has the same formal variance as
the antibracket in Chern--Simons-type master actions.  In quadratic
benchmark models this variance matches the ordinary Chern--Simons
functional; in non-quadratic examples the operations $m_{n-1}$ are
higher cyclic vertices.  An equality of BV theories requires the
Costello--Gwilliam perturbative quantization formalism \cite{CG17}
and a specified local-to-global factorisation model.  This chapter
uses only the algebraic BV operation on the bar complex.
\end{remark}

For a cyclic non-quadratic transferred model, the formal BV action has
the typed expansion
\[
S_{\mathrm{BV,form}}(\alpha)
=\sum_{n\geq1}\frac{1}{(n+1)!}
\left\langle \alpha,\,\ell_n(\alpha,\ldots,\alpha)
\right\rangle_{\mathrm{cyc}}.
\]
The summand with $\ell_3$ is the first place where a cubic
RTT/Serre relation contributes a quartic action term.

\subsection{Implications for Koszul duality}

The Chern--Simons input is external to Koszul duality.  Quadratic
algebras give the ordinary dg~Lie Maurer--Cartan equation,
non-quadratic algebras give higher $L_\infty$ vertices, and the
critical level gives a screening/free-field complex.  Interpreting
algebra/coalgebra exchange through a level-rank-type
Chern--Simons dictionary is a separate conjectural physics comparison,
not a theorem of the bar coalgebra and not a literal identification
with classical level-rank duality.

\begin{remark}[Cumulant duals for nonquadratic algebras]
\label{rem:cumulant-dual-nonquadratic}
\index{cumulant coalgebra!as Koszul dual}
\index{nonquadratic algebra!cumulant dual}
For quadratic chiral algebras ($\Delta_\cA(t) = 0$), the Koszul
dual coalgebra~$\cA^i$ is computed from the quadratic orthogonal
relations, and the bar coalgebra resolves it:
$\barB(\cA)\xrightarrow{\sim}\cA^i$.  The dual algebra
$\cA^!=(\cA^i)^\vee$ is obtained only after the finite-type duality
is available. For genuinely nonquadratic algebras
($\Delta_\cA(t) \neq 0$), the bar coalgebra develops primitive
cogenerators not present among the declared generators of~$\cA$.
The prototype is the $\mathcal{W}_3$ algebra: the composite field
$\Lambda = :(TT): - \tfrac{3}{10}\partial^2 T$ appears as a
primitive cogenerator $\Lambda^*$ in the bar complex, with nonprimitive
coproduct $\Delta(\Lambda^*) = L^* \otimes L^* - \tfrac{3}{10}
(\partial^2 L^* \otimes 1 + 1 \otimes \partial^2 L^*)$ and
differential encoding the characteristic $W$--$W$ OPE coefficient
$16/(22 + 5c)$.

The correct completed dual-coalgebra candidate of a nonquadratic
chiral algebra is therefore the \emph{cumulant coalgebra}
$\operatorname{Cum}_c(\cA) = \hat{T}^c(sQ(\cA)^\vee)$
(Definition~\ref{def:cumulant-coalgebra}), where $Q(\cA) =
\bar{\cA}/\partial\bar{\cA}$ is the canonical primitive sector
extracted from the vacuum character
(Definition~\ref{def:primitive-cumulant-quotient}). The primitive
defect $\Delta_\cA(t)$
(Definition~\ref{def:primitive-defect-series}) measures the
discrepancy between the cumulant coalgebra and the naive dual on
declared generators. The cumulant recognition conjecture
(Conjecture~\ref{conj:cumulant-recognition}) predicts
$\operatorname{gr}_\rho \hat{B}(\cA) \simeq
\operatorname{Cum}_c(\cA)$.

Every algebraic structure is an MC element in
a convolution dg Lie algebra, and the cumulant coalgebra is the
correct \emph{source} coalgebra for those convolution maps.  The
dual algebra is obtained from this coalgebra only after the
Verdier/completed-cobar hypotheses are imposed. Naive generators give
the wrong source; cumulant generators give the canonical one.
\end{remark}

% ================================================================
% THE ORDERED ASSOCIATIVE CORE
% ================================================================

\section{The ordered associative core}
\label{sec:ordered-associative-core}

The chiral Koszul duality of the preceding sections has an ordered
associative core that is invisible from the commutative vantage point.
The geometric origin is simple: the ordered configuration space
$\Conf_n^{\mathrm{ord}}(X) = \{(x_1, \ldots, x_n) \in X^n :
x_i \neq x_j\}$ carries an associative, not merely $E_1$, structure,
because the labelling of points is ordered, not just
unordered-up-to-homotopy. Shuffles survive as distinct summands; they
are not killed by a symmetric-group quotient. The bar coalgebra
$C = \barBch(\cA)$ inherits this ordering, and once the ordering is
taken seriously, a new object appears on the coalgebra side: the
\emph{diagonal bicomodule} $C_\Delta$, which is $C$ equipped with its
coproduct on both left and right. The diagonal bimodule
${}_\cA\cA_\cA$ and the diagonal bicomodule $C_\Delta$ are not the
same object; the two-sided twisting functor
$\mathcal K_\cA^{\mathrm{bi}}$ transports the former to the latter.
Only after the finite-type Verdier comparison identifies the dual
coalgebra lane with~$\cA^!$ may this coalgebraic boundary be rewritten
as an $\cA$--$\cA^!$ interface.

The ordered associative
theory runs from the chiral shuffle theorem through the genus-zero open
trace formalism to the conjectural frontier connecting
it to the modular Maurer--Cartan theory of
Chapter~\ref{chap:higher-genus}.

\begin{remark}[Koszul pairs as transverse boundary conditions]
\label{rem:dimofte-transverse-boundary}
\index{boundary conditions!transverse}%
\index{Dimofte!slab geometry}%
\index{Koszul pair!as transverse boundaries}%
In the holomorphic-topological framework of
Costello--Dimofte--Gaiotto~\cite{CDG2024,CG17}, a
$3$d HT theory on $\bC \times \bR$ admits boundary conditions
at each end of the $\bR$-interval. The intrinsic algebraic datum of a
Koszul pair is the bar coalgebra $C=\barBch(\cA)$ and its diagonal
bicomodule~$C_\Delta$.  The statement that $\cA$ and~$\cA^!$ form
transverse boundary conditions is the finite-type Verdier
translation of this datum, not the definition of the bar construction:
$\cA^!$ enters only after the post-Verdier factorization algebra has
been identified. Bar-cobar
inversion (Theorem~B) is the counit
$\Omegach C\to\cA$ and recovers~$\cA$ itself.  Verdier intertwining
(Theorem~A) relates the bar coalgebra of~$\cA$ to the
post-Verdier algebra lane from which~$\cA^!$ is obtained.
\end{remark}

\subsection{The chiral shuffle theorem and opposite duality}

The first point of departure from the commutative story is an ordered
Eilenberg--Zilber theorem for chiral bar coalgebras.

\begin{theorem}[Ordered chiral shuffle theorem;
\ClaimStatusProvedHere]
\label{thm:ordered-shuffle}
Let $\cA$ and $\mathcal{B}$ be augmented pro-nilpotent associative
chiral algebras on a smooth curve~$X$, with filtrations satisfying
the standard admissibility conditions of
Theorem~\textup{\ref{thm:chiral-bar-resolution-exact}}.
There is a natural filtered quasi-isomorphism of conilpotent dg
coalgebras
\[
 \barBch(\cA) \widehat\otimes \barBch(\mathcal{B})
 \xrightarrow{\;\sim\;}
 \barBch(\cA \chirtensor \mathcal{B}),
\]
realized by shuffle and Alexander--Whitney maps
$\mathsf{Sh}^{\mathrm{ch}}$ and $\mathsf{AW}^{\mathrm{ch}}$
satisfying $\mathsf{AW} \circ \mathsf{Sh} = \mathrm{id}$ and
$\mathsf{Sh} \circ \mathsf{AW} \simeq \mathrm{id}$. Both maps are
coalgebra morphisms.
\end{theorem}

\begin{proof}
Form the bisimplicial object
$K_{\bullet,\bullet} = B^\Delta_\bullet(\cA) \chirtensor
B^\Delta_\bullet(\mathcal{B})$, where $B^\Delta_n(\cA) =
\cA^{\chirtensor n}$ is the simplicial ordered bar object. Its
diagonal satisfies $\operatorname{diag}(K)_n \cong
(\cA \chirtensor \mathcal{B})^{\chirtensor n}$ because
multiplication and augmentation on
$\cA \chirtensor \mathcal{B}$ are componentwise. Apply the
Eilenberg--Zilber theorem filtration-wise in each graded piece;
the normalized-chains functor commutes with the completed chiral
tensor constructions under the admissibility hypotheses. The
resulting shuffle map respects deconcatenation because cutting a
shuffled word is equivalent to cutting the source words and then
shuffling the left and right halves.
\end{proof}

The second foundational phenomenon is order reversal: it transforms
opposite multiplication on the algebra side into opposite
deconcatenation on the bar side.

\begin{theorem}[Opposite duality for ordered bar coalgebras;
\ClaimStatusProvedHere]
\label{thm:ordered-opposite}
For every augmented pro-nilpotent associative chiral algebra~$\cA$,
there is a canonical isomorphism of coaugmented conilpotent dg
coalgebras
\[
 \mathsf{R}_\cA \colon
 \barBch(\cA^{\mathrm{op}})
 \xrightarrow{\;\sim\;}
 \barBch(\cA)^{\mathrm{cop}},
\]
induced by the order-reversal involution
$[sa_1|\cdots|sa_n] \mapsto
(-1)^{\sum_{i<j}(|a_i|+1)(|a_j|+1)}
[sa_n|\cdots|sa_1]$
on suspended bar words. If $\cA$ is finite-type Koszul, then
$(\cA^{\mathrm{op}})^! \simeq (\cA^!)^{\mathrm{op}}$.
\end{theorem}

\begin{proof}
Reversal carries the $i$-th face map of the simplicial bar
object for $\cA^{\mathrm{op}}$ to the $(n-i)$-th face map for $\cA$,
with the stated Koszul sign. The coproduct identity
$\Delta \circ \mathsf{R}_\cA =
\tau \circ (\mathsf{R}_\cA \otimes \mathsf{R}_\cA) \circ \Delta$
(where $\tau$ flips tensor factors) follows because reversal sends a
cut after the $p$-th letter to a cut before the $(n-p)$-th letter.
The statement on Koszul duals follows by finite-type linear duality.
\end{proof}

\begin{corollary}[Enveloping duality; \ClaimStatusProvedHere]
\label{cor:ordered-enveloping}
Set $C = \barBch(\cA)$ and
$C^e = C \widehat\otimes C^{\mathrm{cop}}$. Then
$\barBch(\cA^e) \xrightarrow{\sim} C^e$, and if $\cA$ is
finite-type Koszul, then $(\cA^e)^! \simeq (\cA^!)^e$.
\end{corollary}

\begin{proof}
Combine the ordered shuffle theorem
(Theorem~\ref{thm:ordered-shuffle}) with opposite duality
(Theorem~\ref{thm:ordered-opposite}).
\end{proof}

\subsection{The diagonal bicomodule and the bimodule--bicomodule
equivalence}

The diagonal bicomodule $C_\Delta$ is $C$ equipped with left and
right coactions both equal to the deconcatenation coproduct
$\Delta_C$. It is the coalgebra-side avatar of the diagonal bimodule
${}_\cA \cA_\cA$.

\begin{definition}[Two-sided twisted bicomodule]
\label{def:two-sided-twisted-bicomodule}
Let $\cM$ be an $\cA$-bimodule
(Definition~\ref{def:chiral-bimodule}). The \emph{two-sided twisted
bicomodule} is
\[
 \mathcal{K}_\cA^{\mathrm{bi}}(\cM)
 = C \widehat\otimes \cM \widehat\otimes C,
\]
with differential twisted by the universal twisting cochain
$\tau \colon C \to \cA$ via left and right action terms. The left
and right $C$-coactions are deconcatenation on the two outer factors.
The Maurer--Cartan equation $d\tau + \tau \star \tau = 0$ ensures
$d^2 = 0$ on the uncurved or strictly corrected completed lane.  In
the genus-$g$ curved lane the same construction is read as a CDG
bicomodule with square governed by the central curvature term, before
passing to the coderived comparison.
\end{definition}

\begin{theorem}[Bimodule--bicomodule equivalence;
\ClaimStatusConditional]
\label{thm:ordered-bimod-bicomod}
Assume that $\cA$ is augmented and pro-nilpotent, that the ordered
shuffle map of Theorem~\textup{\ref{thm:ordered-shuffle}} applies to
$\cA^e$, and that the PBW-complete bar coalgebra
$C^e=\barBch(\cA^e)$ satisfies the conilpotence and
co/contra exactness hypotheses of
Theorem~\textup{\ref{thm:positselski-chiral}}.  Then there is an
exact equivalence
\[
 \mathcal{K}_\cA^{\mathrm{bi}} \colon
 D^{\mathrm{pbw}}_{\mathrm{bi}}(\cA)
 \xrightarrow{\;\sim\;}
 D^{\mathrm{co}}_{\mathrm{bi}}(C)
\]
between the derived category of PBW-complete $\cA$-bimodules and
the coderived category of conilpotent $C$-bicomodules, with inverse
given by the two-sided cobar functor
$\cM \mapsto \cA \widehat\otimes \cM \widehat\otimes \cA$.
\end{theorem}

\begin{proof}
By Corollary~\ref{cor:ordered-enveloping},
$\barBch(\cA^e) \simeq C^e$.
Apply the one-sided complete bar-coalgebra comparison to the
enveloping algebra $\cA^e$.  The target is the coderived category of
comodules over the bar coalgebra $C^e$, not an ordinary module
category over~$\cA^!$ unless the separate finite-type dualization
identifies $C^*$ with~$\cA^!$, equivalently identifies
$(C^e)^*$ with~$(\cA^!)^e$. Since left $C^e$-comodules are
canonically $C$-bicomodules, the result follows.
\end{proof}

\begin{theorem}[Diagonal correspondence; \ClaimStatusConditional]
\label{thm:ordered-diagonal}
Under the equivalence of
Theorem~\textup{\ref{thm:ordered-bimod-bicomod}}, the diagonal
bimodule ${}_\cA \cA_\cA$ maps to the diagonal bicomodule
$C_\Delta$. Equivalently,
$\mathcal{K}_\cA^{\mathrm{bi}}(\cA) \simeq C_\Delta$ in
$D^{\mathrm{co}}_{\mathrm{bi}}(C)$.
\end{theorem}

\begin{proof}
The bar--cobar counit contractions
$C \widehat\otimes_\tau \cA \xrightarrow{\sim} \mathbf{1}$ and
$\cA \widehat\otimes_\tau C \xrightarrow{\sim} \mathbf{1}$
collapse the two-sided twisted bicomodule
$C \widehat\otimes \cA \widehat\otimes C$ to a single copy of~$C$.
Filter by total coradical degree in the two outer copies; on the
associated graded, the twisting terms vanish and both coactions are
the two legs of~$\Delta_C$. Because the filtration is complete and
bounded below, this identification lifts to the filtered complex.
\end{proof}

The diagonal bicomodule is the composition unit for derived
cotensor: for every $X \in D^{\mathrm{co}}_{\mathrm{bi}}(C)$,
\[
 C_\Delta \mathbin{\square}_C^{\mathbf{R}} X
 \simeq X
 \simeq X \mathbin{\square}_C^{\mathbf{R}} C_\Delta.
\]
This is the transport of the tautology
$\cA \otimes_\cA^{\mathbf{L}} \cM \simeq \cM$ across the
equivalence. Tensor--cotensor
gluing $\mathcal{K}_\cA^{\mathrm{bi}}(\cM \otimes_\cA^{\mathbf{L}}
\mathcal{N}) \simeq \mathcal{K}_\cA^{\mathrm{bi}}(\cM)
\mathbin{\square}_C^{\mathbf{R}}
\mathcal{K}_\cA^{\mathrm{bi}}(\mathcal{N})$ makes $C_\Delta$ the
universal boundary condition for interval composition.

\subsection{chiral Hochschild--coHochschild duality}

The chiral Hochschild cohomology
$\ChirHoch^*(\cA, \cM) = \Ext^*_{\cA^e}(\cA, \cM)$ constructed in
Theorem~\ref{thm:chiral-hochschild-complex}, and computed on the bar
resolution in Theorem~\ref{thm:chiral-hochschild-complex}, has an
exact coalgebraic counterpart: the coHochschild theory of the bar
coalgebra.

\begin{theorem}[chiral Hochschild--coHochschild duality, homological
version; \ClaimStatusConditional]
\label{thm:ordered-HH-coHH-homology}
For every complete $\cA$-bimodule~$\cM$,
\[
 \mathrm{HH}^{\mathrm{ch}}_\bullet(\cA, \cM)
 \;\simeq\;
 \mathrm{coHH}^{\mathrm{ch}}_\bullet
 \bigl(C, \mathcal{K}_\cA^{\mathrm{bi}}(\cM)\bigr),
\]
where
$\mathrm{coHH}^{\mathrm{ch}}_\bullet(C, N) =
\mathrm{Cotor}_\bullet^{C^e}(C_\Delta, N)$.
In particular,
$\mathrm{HH}^{\mathrm{ch}}_\bullet(\cA) \simeq
C_\Delta \mathbin{\square}_{C^e}^{\mathbf{R}} C_\Delta$:
the annulus is the self-cotensor of the universal boundary
condition.
\end{theorem}

\begin{proof}
Transport $\mathrm{Tor}_\bullet^{\cA^e}(\cA, \cM)$ across the
bimodule--bicomodule equivalence
(Theorem~\ref{thm:ordered-bimod-bicomod}). The diagonal
bimodule~$\cA$ maps to~$C_\Delta$ by
Theorem~\ref{thm:ordered-diagonal}; $\cM$ maps to
$\mathcal{K}_\cA^{\mathrm{bi}}(\cM)$ by definition.
Derived tensor over~$\cA^e$ becomes derived cotensor
over~$C^e$.
\end{proof}

\begin{theorem}[chiral Hochschild--coHochschild duality, cohomological
version; \ClaimStatusProvedHere]
\label{thm:ordered-HH-coHH-cohomology}
Under the same hypotheses,
\[
 \ChirHoch^*(\cA, \cM)
 \;\simeq\;
 \mathrm{Coext}_{C^e}^*
 \bigl(C_\Delta,
 \mathcal{K}_\cA^{\mathrm{bi}}(\cM)\bigr).
\]
In particular,
$\ChirHoch^*(\cA) \simeq
\mathrm{Coext}_{C^e}^*(C_\Delta, C_\Delta)$: the chiral
Hochschild cohomology ring whose Gerstenhaber bracket was
constructed in Theorem~\textup{\ref{thm:chiral-gerstenhaber-kps}} is
the coalgebraic self-extension of the diagonal.
\end{theorem}

\begin{proof}
Derived Hom is preserved by equivalences of derived categories.
Apply the bimodule--bicomodule equivalence to
$R\!\mathrm{Hom}_{\cA^e}(\cA, \cM)$.
\end{proof}

\begin{remark}[Relation to existing chiral Hochschild material]
\label{rem:ordered-HH-relation}
The cohomological version recovers the complex of
Theorem~\ref{thm:chiral-hochschild-complex} as the coalgebraic
coHochschild complex of~$C$. The Gerstenhaber cup product of
Proposition~\ref{prop:cup-product-properties} is the composition
product on $\mathrm{Coext}_{C^e}^*(C_\Delta, C_\Delta)$, and the
periodicity theorems
(Theorem~\ref{thm:virasoro-periodicity},
Theorem~\ref{thm:affine-periodicity-critical}) become statements
about the coalgebraic self-extension ring. The exchange
$\ChirHoch^*(\cA) \leftrightarrow \ChirHoch^*(\cA^!)$ of
Corollary~\ref{cor:hochschild-cup-exchange} is the finite-type
Verdier shadow of enveloping duality
(Corollary~\ref{cor:ordered-enveloping}); it is not a consequence of
the bar-cobar counit alone.
\end{remark}

\subsection{The pair-of-pants product and the open trace formalism}

The diagonal bicomodule carries a second algebraic operation beyond
cotensor composition: the \emph{ordered pair-of-pants product},
transported from the multiplication
$\mu_\cA \colon \cA \chirtensor \cA \to \cA$.

\begin{theorem}[Ordered pair-of-pants algebra;
\ClaimStatusConditional]
\label{thm:ordered-pair-of-pants}
There exist an associative multiplication
$\mu_P \colon C_\Delta \star C_\Delta \to C_\Delta$
and a unit
$\eta_P \colon \mathbf{1}_\star \to C_\Delta$
for the ordered fusion product
$X \star Y = \mathcal{K}_\cA^{\mathrm{bi}}(
\Omega_C^{\mathrm{bi}}(X) \widehat\otimes
\Omega_C^{\mathrm{bi}}(Y))$,
obtained by transporting $(\mu_\cA, \eta_\cA)$ across
$\mathcal{K}_\cA^{\mathrm{bi}}$ and the diagonal correspondence.
\end{theorem}

\begin{proof}
The diagonal bimodule ${}_\cA \cA_\cA$ is an associative algebra
object for the tensor product of bimodules. Apply the exact
equivalence of
Theorem~\ref{thm:ordered-bimod-bicomod}: associativity and
unitality are preserved because the equivalence preserves
composition.
\end{proof}

\begin{theorem}[Master theorem: the ordered open trace formalism;
\ClaimStatusConditional]
\label{thm:ordered-master}
Let $\cA$ be an augmented pro-nilpotent associative chiral algebra
on~$X$ satisfying the admissibility conditions, and set
$C = \barBch(\cA)$. Then the data
\[
 \bigl(D^{\mathrm{co}}_{\mathrm{bi}}(C),\;
 \mathbin{\square}_C^{\mathbf{R}},\;
 C_\Delta,\;
 \star,\;
 \mu_P,\;
 \eta_P,\;
 \mathcal{H}_C^{\mathrm{cl}}\bigr)
\]
constitute an ordered genus-zero open trace formalism:
\begin{enumerate}
\item $C_\Delta$ is the unit for interval composition
 $\mathbin{\square}_C^{\mathbf{R}}$.
\item $(C_\Delta, \mu_P, \eta_P)$ is an associative algebra object
 for ordered fusion~$\star$.
\item The closed object is the coHochschild complex:
 $\mathcal{H}_C^{\mathrm{cl}} \simeq
 \mathrm{coHH}^{\mathrm{ch}}_\bullet(C) \simeq
 \mathrm{HH}^{\mathrm{ch}}_\bullet(\cA)$.
\item The annulus acts on $C_\Delta$ via the cap action transported
 from the chiral Hochschild cap on~$\cA$.
\item All genus-zero sewing identities on the coalgebra side are
 transported from the corresponding identities for
 $\cA$-bimodules.
\end{enumerate}
If $\cA$ is $d$-Calabi--Yau \textup(i.e.\ perfect as an
$\cA^e$-module with a shifted self-duality
$\cA \xrightarrow{\sim} R\!\mathrm{Hom}_{\cA^e}(\cA, \cA^e)[d]$\textup),
then the underlying complex of $C_\Delta$ acquires a $d$-shifted
ordered Frobenius structure: the trace $\varepsilon_P \colon C_\Delta \to
k[-d]$ is transported from the Calabi--Yau trace, the coproduct
$\Delta_P$ is adjoint to $\mu_P$ under $\varepsilon_P$, and the
Frobenius identity $(\mu_P \star \mathrm{id}) \circ
(\mathrm{id} \star \Delta_P) = \Delta_P \circ \mu_P$ holds.
\end{theorem}

\begin{proof}
This compresses
Theorems~\ref{thm:ordered-shuffle}--\ref{thm:ordered-pair-of-pants}
together with the Hochschild--coHochschild dualities
(Theorems~\ref{thm:ordered-HH-coHH-homology}
and~\ref{thm:ordered-HH-coHH-cohomology}). The Calabi--Yau
refinement: transport the algebra $(\mu_\cA, \eta_\cA)$ and its
adjoint coalgebra $(\Delta_\cA, \varepsilon_\cA)$ simultaneously;
the Frobenius identities on $C_\Delta$ are the images of the
Frobenius identities on~$\cA$, which are formal consequences of
associativity and nondegeneracy.
\end{proof}

\begin{remark}[The dictionary]
\label{rem:ordered-dictionary}
The ordered associative core provides the following dictionary
between the algebra~$\cA$ and the coalgebra~$C$:
\[
 \cA \longleftrightarrow C_\Delta,
 \qquad
 \cA^e \longleftrightarrow C^e,
 \qquad
 {-}\otimes_\cA^{\mathbf{L}}{-}
 \longleftrightarrow
 {-}\mathbin{\square}_C^{\mathbf{R}}{-},
\]
\[
 \mathrm{HH}^{\mathrm{ch}}_\bullet(\cA, \cM)
 \longleftrightarrow
 \mathrm{coHH}^{\mathrm{ch}}_\bullet
 (C, \mathcal{K}_\cA^{\mathrm{bi}}(\cM)),
 \qquad
\ChirHoch^*(\cA, \cM)
 \longleftrightarrow
 \mathrm{Coext}_{C^e}^*
 (C_\Delta, \mathcal{K}_\cA^{\mathrm{bi}}(\cM)).
\]
Thus the bar coalgebra is the ordered coalgebraic state object, the
diagonal bicomodule is the interval unit, and the annulus is the
self-cotensor of that unit.
\end{remark}

\begin{proposition}[Open trace is not a Drinfeld double;
\ClaimStatusConditional]
\label{prop:open-trace-not-drinfeld-double}
\index{Drinfeld double!separated from open trace}
\index{open trace!not Drinfeld double}
The ordered open trace formalism of
Theorem~\textup{\ref{thm:ordered-master}} does not construct a
Drinfeld double.  For a finite-dimensional Hopf algebra~$H$, the
Drinfeld double is the Hopf algebra
\[
 D(H)=H^{*\operatorname{cop}}\bowtie H
\]
with multiplication determined by the Hopf pairing and the antipode,
and $\operatorname{Rep}D(H)\simeq Z(\operatorname{Rep}H)$.
The closed object in the ordered open trace formalism is instead
\[
 \mathcal H_C^{\mathrm{cl}}
 \simeq \mathrm{coHH}^{\mathrm{ch}}_\bullet(C)
 \simeq C_\Delta\square_{C^e}^{\mathbf R}C_\Delta.
\]
It is a Hochschild/coHochschild trace complex, not a Hopf algebra.
Thus an external Hall--Drinfeld double may act on, or be compared
with, this trace only after a separate Hopf/Hall bialgebra structure
and a compatible pairing have been supplied.
\end{proposition}

\begin{proof}
The open trace data are typed by bicomodules over the bar coalgebra
$C=\barBch(\cA)$ and by derived cotensor over $C^e$.  They require no
multiplication on~$C$, no antipode, and no Hopf pairing
$C^*\otimes C\to k$.  A Drinfeld double requires exactly those Hopf
inputs.  The self-cotensor
$C_\Delta\square_{C^e}^{\mathbf R}C_\Delta$ computes a trace object
in the bicomodule category; its product is the Hochschild cup/cap
structure transported by Theorems~\ref{thm:ordered-HH-coHH-homology}
and~\ref{thm:ordered-HH-coHH-cohomology}, not the crossed
multiplication of $H^{*\operatorname{cop}}\bowtie H$.
\end{proof}

\subsection{The commutator-shadow theorem and the conjectural frontier}

The preceding subsections supply the genus-zero ordered algebraic
core.  The commutator filtration compares this ordered associative
bar coalgebra with the Francis--Gaitsgory Lie/cocommutative bar
coalgebra on the associated graded.  Three independent obstructions
remain outside that comparison.

\begin{theorem}[Commutator-shadow theorem; \ClaimStatusProvedElsewhere]
\label{thm:ordered-FG-shadow}
Let $\cA$ be a filtered strongly admissible associative chiral algebra
with exhaustive separated commutator filtration
$F_{\mathrm{com}}^\bullet \cA$. Assume:
\begin{enumerate}[label=\textup{(\alph*)}]
\item $\operatorname{gr}_{\mathrm{com}} \cA$ is $E_\infty$-chiral;
\item $\operatorname{gr}_{\mathrm{com}} \cA$ is quadratic Koszul on the
 commutative/Lie side;
\item the induced filtration on $\barBch(\cA)$ is complete and
 strongly convergent.
\end{enumerate}
Then there is a convergent spectral sequence
\[
 E_1 \cong \barB^{FG}\!\bigl(\operatorname{gr}_{\mathrm{com}} \cA\bigr)
 \;\Longrightarrow\;
 \operatorname{gr}_{\mathrm{com}} \barBch(\cA),
\]
recovering the Francis--Gaitsgory Lie/cocommutative duality as the
associated graded of the ordered associative theory.
\end{theorem}

\begin{proof}[Proof sketch \textup{(}full proof:
Vol~\textup{II,} Part~\textup{IX;} cf.\ Chriss--Ginzburg
\cite{CG97}\textup{)}]
The argument adapts the Chriss--Ginzburg stratification spectral
sequence from the Steinberg variety to the chiral correspondence
space $\mathrm{FM}_k(\mathbb{C})\times\mathrm{Conf}_k(\mathbb{R})$
underlying the bar complex. Five ingredients combine.

\smallskip
\emph{(a)~Commutator stratification.}
The commutator filtration $F_{\mathrm{com}}^\bullet \cA$ induces
an increasing filtration $S_0\subset S_1\subset\cdots$ of the
correspondence space by commutator-depth strata. For the
associated graded $\operatorname{gr}_{\mathrm{com}}\cA$ (which is
$E_\infty$-chiral), only the lowest stratum~$S_1$ contributes:
all OPE singularities on~$\operatorname{gr}_{\mathrm{com}}\cA$
are at worst simple poles.

\smallskip
\emph{(b)~$E_1$ identification.}
The stratification spectral sequence has $E_1$~page
\[
E_1^{p,q}
\;\cong\;
H^{p+q}_{\mathrm{BM}}(S_p/S_{p-1})
\;\cong\;
\barB^{FG}_{p+q}\!\bigl(\operatorname{gr}_{\mathrm{com}}\cA\bigr).
\]
On the associated graded, only $p=1$ is nontrivial (all higher strata
are empty), so the spectral sequence collapses to the
Francis--Gaitsgory bar complex. The surviving differential on~$E_1$
is the boundary map of the Borel--Moore complex of~$S_1$, which is
the FG differential by construction.

\smallskip
\emph{(c)~Pole non-increase: the commutator filtration.}
The bar differential $d_{\mathrm{bar}}$ respects the commutator
filtration on~$\barBch(\cA)$: the $n=0$ face (Lie bracket) maps
$F^p_{\mathrm{com}}\to F^{p+1}_{\mathrm{com}}$, while the $n\ge 1$
faces preserve $F^p_{\mathrm{com}}$. The proof uses the Borcherds
identity at parameter $m=0$ and the strong-generation bound on
iterated pole orders. In particular, $d_{\mathrm{bar}}(F^p)\subset F^p$
for the total filtration.

\smallskip
\emph{(d)~The Virasoro subtlety.}
For generators of weight~$\ge 2$ (Virasoro, $\mathcal{W}$-algebras),
the \emph{raw} pole-order filtration is not respected by the bar
differential: $T_{(0)}\,T=\partial T$ has $N(\partial T,T)=N(T,T)+1$
(differentiation raises pole order). However, the commutator (PBW)
filtration \emph{is} respected, because the PBW convention treats
$\partial$ as a structural operator at the same filtration level.
The theorem is stated in terms of $F_{\mathrm{com}}^\bullet$, so the
commutator-filtration argument suffices. For the affine lineage
($\hat\fg_k$, free fields, lattice algebras: all generators of
weight~$1$), the two filtrations agree.

\smallskip
\emph{(e)~Convergence.}
Strong convergence of the filtration (hypothesis~(c)) and the
$E_1$~identification yield
$E_1\cong\barB^{FG}(\operatorname{gr}_{\mathrm{com}}\cA)\Rightarrow
\operatorname{gr}_{\mathrm{com}}\barBch(\cA)$, as claimed.
\end{proof}

\begin{corollary}[Verdier-Koszul dual on the associated graded;
\ClaimStatusConditional]
\label{cor:ordered-FG-koszul-locus}
On the finite-type Verdier-Koszul locus, where
$\Bbarch_X(\cA)$ resolves the dual coalgebra~$\cA^i$ and
$\cA^!=(\cA^i)^\vee$ exists,
\[
 \operatorname{gr}_{\mathrm{com}}(\cA^!)
 \;\simeq\;
 \bigl(\operatorname{gr}_{\mathrm{com}} \cA\bigr)^!_{FG}.
\]
That is, the Koszul dual of the associated graded (computed via
Francis--Gaitsgory Lie/cocommutative duality) recovers the
associated graded of the Koszul dual.
\end{corollary}

\begin{proof}
On this locus, Koszulness gives a quasi-isomorphism of coalgebras
$\barBch(\cA)\xrightarrow{\sim}\cA^i$, and the finite-type
Verdier/linear-dual comparison gives
$\cA^!=(\cA^i)^\vee\simeq\barBch(\cA)^\vee$.  This use of the bar
coalgebra is separate from the bar-cobar counit
$\Omega\barBch(\cA)\to\cA$, whose target is~$\cA$.
The filtration $F_{\mathrm{com}}^\bullet$ is exhaustive and
separated (hypothesis~(c) of
Theorem~\ref{thm:ordered-FG-shadow}), so the associated graded of the
finite-type dual is the dual of the associated graded:
$\operatorname{gr}_{\mathrm{com}}(\cA^!)
\simeq\operatorname{gr}_{\mathrm{com}}(\barBch(\cA))^\vee
\simeq\barB^{FG}(\operatorname{gr}_{\mathrm{com}}\cA)^\vee
=(\operatorname{gr}_{\mathrm{com}}\cA)^!_{FG}$.
\end{proof}

\begin{conjecture}[Bordered realization; \ClaimStatusConjectured]
\label{conj:ordered-bordered}
There exists a bordered ordered configuration-space model
$\mathfrak{C}^{\mathrm{bord}}_{g,r,s}(X)$ whose chain-level
operations recover the ordered open trace formalism of
Theorem~\textup{\ref{thm:ordered-master}}: interval components give
$C_\Delta$, boundary gluing is derived cotensor, the pair-of-pants
gives $\mu_P$, and annular closure gives
$\mathrm{coHH}^{\mathrm{ch}}_\bullet(C)$.
\end{conjecture}

\begin{conjecture}[Associative modular Maurer--Cartan class;
\ClaimStatusConjectured]
\label{conj:ordered-modular-MC}
There exists an associative cyclic deformation complex
$\mathrm{Def}^{\mathrm{cyc}}_{\Ass^{\mathrm{ch}}}(\cA)$
and a Maurer--Cartan element
$\Theta_\cA^{\Ass^{\mathrm{ch}}}$ whose linear term is the
genus anomaly, whose quadratic term is the first pair-of-pants
boundary correction, and whose full Maurer--Cartan equation
expresses compatibility of all ordered higher-genus boundary
gluings.  The class is required to restrict to the
ordered-associative projection of the modular MC class~$\Theta_\cA$
of Chapter~\textup{\ref{chap:higher-genus}}.
\end{conjecture}

\begin{conjecture}[Drinfeld--Sokolov compatibility;
\ClaimStatusConjectured]
\label{conj:ordered-DS}
BRST/Drinfeld--Sokolov reduction commutes with ordered associative
chiral duality: $Q_{DS}(\cA^!) \simeq Q_{DS}(\cA)^!$, and the
two-sided twisted bicomodule functor
$\mathcal{K}_\cA^{\mathrm{bi}}$ intertwines reduction on bimodules
with reduction on bicomodules.
\end{conjecture}

Theorem~\ref{thm:ordered-FG-shadow} and
Corollary~\ref{cor:ordered-FG-koszul-locus} identify
Francis--Gaitsgory duality as the associated graded of the ordered
theory under the commutator filtration.  The remaining layer has
three independent obstructions.  Conjecture~\ref{conj:ordered-bordered} asks for a
bordered configuration-space substrate.  Conjecture~\ref{conj:ordered-modular-MC}
asks for a higher-genus ordered MC class refining the shadow
obstruction tower of Chapter~\ref{chap:higher-genus}.  Conjecture~\ref{conj:ordered-DS}
asks for BRST/Drinfeld--Sokolov compatibility; if reduction commutes
with duality, the ordered open trace formalism descends to the
reduced theory.

% ================================================================
% SECTION: CATEGORICAL KOSZULNESS AND THE MODULE-LEVEL FRONTIER
% ================================================================

\section{Categorical Koszulness and the module-level frontier}
\label{sec:categorical-koszulness-frontier}

The module-level comparison statements of
Theorems~\ref{thm:koszul-equivalence-categories}--\ref{thm:full-derived-module-equiv}
are triangulated. The $\infty$-categorical
refinement comprises t-exactness of bar duality, weight decomposition of module
categories, the chiral BGG correspondence, algebraicity of bar generating
functions, and the interaction of Koszulness with $C_2$-cofiniteness.

\subsection{Module category t-exactness}

\begin{conjecture}[Module category t-exactness; \ClaimStatusConjectured]
\label{conj:module-category-t-exactness}
\index{t-exactness!module category}
\index{Koszul pair!module category t-exactness}
Let $(\cA,\cA^i)$ be a chiral Koszul pair on a smooth curve~$X$.
Set $C=\barB(\cA)$.  The bar functor
$\Phi(M) = C \otimes_{\cA} M$ induces
an equivalence of stable $\infty$-categories
\[
 \Phi \colon \mathrm{FactMod}(\cA)
 \xrightarrow{\;\sim\;}
 \mathrm{FactCoMod}(C).
\]
If the finite-type Verdier comparison identifies
$C\simeq\cA^i=(\cA^!)^\vee$, the target may be rewritten as a
conilpotent $\cA^i$-comodule category, or dually as a completed
$\cA^!$-module category. Equip $\mathrm{FactMod}(\cA)$ with the
\emph{PBW t-structure}
and $\mathrm{FactCoMod}(C)$ with the \emph{cobar t-structure}
\textup{(}defined below\textup{)}. Then $\cA$ is chirally Koszul
if and only if $\Phi$ is t-exact.
\end{conjecture}

\textbf{The two t-structures.}
The PBW filtration $F_\bullet$ on $\mathrm{FactMod}(\cA)$ is the
decreasing filtration by bar degree: $F_p M$ is the submodule
generated by elements of bar filtration $\geq p$ (i.e., those
requiring $\geq p$ iterated OPE contractions to reach the vacuum).
A module $M$ lies in $\mathcal{D}_{\mathrm{PBW}}^{\leq 0}$ if
its associated graded $\mathrm{gr}^F_p M$ is acyclic for $p > 0$.
The cobar t-structure on $\mathrm{FactCoMod}(C)$ is defined
dually: a comodule $N$ lies in $\mathcal{D}_{\mathrm{cob}}^{\leq 0}$
if $\mathrm{gr}^G_p N$ is acyclic for $p > 0$, where $G_\bullet$
is the cobar filtration.

\smallskip
\noindent\emph{Equivalences.}
\begin{enumerate}[leftmargin=2em]
\item t-exactness of $\Phi$ $\iff$ the PBW spectral sequence
 collapses at $E_2$ for \emph{all} factorization modules~$M$.
\item Heart preservation $\iff$ Koszul linearity:
 $\barB_n(L(\lambda))$ is concentrated in internal degree
 $n + h(\lambda)$, where $h(\lambda)$ is the conformal weight.
\end{enumerate}

\noindent
The conditional highest-weight shadow is
Proposition~\ref{conj:koszul-t-structures}, which records the
transported Koszul t-structure once the conjectural module-category
$t$-exactness and linearity input are supplied. The $\infty$-categorical
upgrade uses the model structure from
Theorem~\ref{thm:quillen-equivalence-chiral}; the
Ch$_\infty$-rectification
(Corollary~\ref{cor:rectification-ch-infty}) ensures that
homotopy and strict module categories agree.
The no-bifunctor obstruction (RNW19~\S6) forces the construction
of $\Phi$ to proceed one slot at a time.

\subsection{Semi-orthogonal decomposition by weight}

\begin{conjecture}[Semi-orthogonal decomposition by conformal weight; \ClaimStatusConjectured]
\label{conj:sod-by-weight}
\index{semi-orthogonal decomposition!by conformal weight}
\index{Koszul pair!weight decomposition}
For a chirally Koszul algebra~$\cA$, the stable $\infty$-category
$\mathrm{FactMod}(\cA)$ admits a semi-orthogonal decomposition
\textup{(}in the sense of Bondal--Orlov\textup{)}
\[
 \mathrm{FactMod}(\cA) = \langle \mathcal{C}_0,\, \mathcal{C}_1,\,
 \mathcal{C}_2,\, \ldots \rangle
\]
where $\mathcal{C}_n$ is the full subcategory of modules generated in
conformal weight~$n$, and the semi-orthogonality condition is
$\mathrm{Hom}(M, N) = 0$ for $M \in \mathcal{C}_i$,
$N \in \mathcal{C}_j$, $i > j$. The decomposition is \emph{split}:
the inclusion $\mathcal{C}_n \hookrightarrow \mathrm{FactMod}(\cA)$
admits exact left and right adjoints.
\end{conjecture}

\noindent
The weight filtration $\mathcal{F}_n = \{M \mid M$
generated in weights $\leq n\}$ has graded pieces
$\mathcal{C}_n = \mathcal{F}_n / \mathcal{F}_{n-1}$.
Koszulness is equivalent to vanishing of cross-weight morphisms:
bar linearity forces the bar differential to preserve weight strata,
killing all $\mathrm{Ext}^*(\mathcal{C}_i, \mathcal{C}_j)$ for
$i > j$. Off the Koszul locus, nonlinear bar terms source
cross-weight extensions and the SOD fails.

\subsection{The chiral BGG correspondence}

\begin{conjecture}[Chiral BGG correspondence; \ClaimStatusConjectured]
\label{conj:chiral-bgg}
\index{BGG correspondence!chiral}
\index{affine Grassmannian!BGG functor}
For a chiral Koszul pair $(\cA,\cA^i)$ with
$\cA = \AffKM{g}_k$ and finite-type Verdier dual~$\cA^!$ when
available,
there is an exact functor
\[
 \mathrm{BGG} \colon \mathcal{O}_k
 \longrightarrow D^b(\mathrm{Coh}(\mathrm{Gr}_G))
\]
from category~$\mathcal{O}$ at level~$k$ to the derived category
of coherent sheaves on the affine Grassmannian, factoring as
\[
 \mathcal{O}_k
 \xrightarrow{\;\barB\;}
 \mathrm{CoMod}^{\mathrm{conil}}(C)
 \xrightarrow{\;\mathrm{Sat}\;}
 D^b(\mathrm{Coh}(\mathrm{Gr}_G)),
\]
where $C=\barB(\cA)$, the first arrow is the bar resolution, and the
second arrow is available only after the separate finite-type
Verdier/Satake comparison identifies the relevant conilpotent
bar-coalgebra comodules with sheaves on~$\mathrm{Gr}_G$.
\begin{enumerate}[label=\textup{(\roman*)},leftmargin=2.5em]
\item Verma modules $M(\lambda) \mapsto \mathcal{O}_{S_\lambda}$
 \textup{(}structure sheaves of Schubert cells\textup{)}.
\item The BGG resolution of $L(\lambda)$ maps to the Cousin
 resolution of $\mathrm{IC}(\overline{S}_\lambda)$.
\item Koszul linearity forces the Schubert filtration spectral
 sequence on $\mathrm{IC}(\overline{S}_\lambda)$ to collapse
 at~$E_1$; the $E_1$ page \emph{is} the bar resolution.
\end{enumerate}
\end{conjecture}

\begin{remark}[Bar side and factorization-Satake obligation]
\label{rem:chiral-bgg-proof-obligation}
The bar side of Conjecture~\ref{conj:chiral-bgg} is established:
Theorem~\ref{thm:bgg-from-bar} constructs the BGG resolution from
the bar complex, Proposition~\ref{prop:bar-bgg-sl2} gives the
explicit bar--BGG for $\AffKM{sl}_{2,k}$, and
Corollary~\ref{cor:bgg-koszul-involution} identifies the BGG
involution under Koszul duality.

The remaining open mathematical condition is the second arrow: the geometric
Satake equivalence must be lifted to the
\emph{factorization} level (identifying conilpotent comodules over the
bar-dual coalgebra, and only after finite-type dualization the
$\cA^!$-side, with constructible sheaves on~$\mathrm{Gr}_G$). This
factorization-Satake lift is not constructed here. The conjecture is
therefore a precise target for connecting chiral Koszul duality to
geometric Langlands.
\end{remark}

\subsection{Bar generating function algebraicity}

\begin{conjecture}[Bar generating function algebraicity; \ClaimStatusConjectured]
\label{conj:bar-gf-algebraicity}
\index{bar generating function!algebraicity}
\index{discriminant!bar generating function}
A chiral algebra~$\cA$ is chirally Koszul if and only if its
bar cohomology generating function
$P_\cA(x) = \sum_{n \geq 0} \dim H^n(\barB(\cA)) \cdot x^n$
satisfies a polynomial equation of degree $\leq 2$:
\begin{equation}\label{eq:bar-gf-quadratic}
 c_2(x) \, P_\cA^2 + c_1(x) \, P_\cA + c_0(x) = 0
\end{equation}
with discriminant
$\Delta_\cA(x) = c_1(x)^2 - 4\, c_0(x)\, c_2(x)$
a polynomial of degree $\leq 2 \cdot \mathrm{rank}(\cA)$.
\end{conjecture}

\begin{proposition}[Computed discriminant rows and the first
non-simply-laced obligation; \ClaimStatusConditional]
\label{prop:bar-gf-discriminant-witnesses-kps}
On the standard computed rows recorded in
Table~\textup{\ref{tab:bar-cohomology-gfs}} and in
Theorem~\textup{\ref{thm:ds-bar-gf-discriminant}}, the bar
generating functions have the following algebraic degrees and
discriminants:
\begin{center}
\renewcommand{\arraystretch}{1.3}
\begin{tabular}{lcc}
\toprule
Family & Algebraic degree & Discriminant $\Delta(x)$ \\
\midrule
$\AffKM{sl}_2$ & 2 & $(1 - 3x)(1 + x)$ \\
Virasoro & 2 & $(1 - 3x)(1 + x)$ \\
$\beta\gamma$ & 2 & $(1 - 3x)(1 + x)$ \\
$\AffKM{sl}_3$ & 2 & $(1 - 8x)(1 - 3x - x^2)$ \\
\bottomrule
\end{tabular}
\end{center}
The shared discriminant $(1 - 3x)(1 + x)$ within the first three
families is the rank-one DS discriminant theorem:
Theorem~\ref{thm:ds-bar-gf-discriminant} proves that
Drinfeld--Sokolov reduction preserves the discriminant, and
$\AffKM{sl}_2$, Virasoro, and $\beta\gamma$ all lie in the same DS
family.

For $\widehat{G}_2$, only $H^1 = 14$ is computed
(Proposition~\ref{prop:G2-bar-dims}); bar cohomology beyond
degree~$1$ is the first open mathematical condition requiring explicit
$G_2$ structure constants.
The conjecture predicts: $P_{\widehat{G}_2}(x)$ satisfies a
quadratic equation with discriminant of degree $\leq 4$ (since
$\mathrm{rank}(G_2) = 2$) and a simple root at $x = 1/14$
(Conjecture~\ref{conj:non-simply-laced-discriminant}).
Verification at $H^2$ would already constrain the quadratic
coefficients.
\end{proposition}

\begin{proof}
The table is the finite witness surface of
Table~\ref{tab:bar-cohomology-gfs}.  The rank-one equality of
discriminants is Theorem~\ref{thm:ds-bar-gf-discriminant}.  The
$\widehat{\mathfrak{sl}}_3$ row is listed there as conditional, so
the proposition keeps the same conditional status.  The final
$G_2$ sentence is not a computed row; it is the next explicit
calculation required to test the conjecture.
\end{proof}

\begin{conjecture}[Chiral Koszul functional equation; \ClaimStatusConjectured]
\label{conj:chiral-koszul-functional-equation}
\index{Koszul functional equation!chiral}
If $(\cA,\cA^i)$ is chirally Koszul and the finite-type
Verdier-dual algebra~$\cA^!$ exists, then the bar generating
functions satisfy the \emph{chiral Koszul functional equation}:
\begin{equation}\label{eq:chiral-koszul-fe}
 P_\cA(x) \cdot P_{\cA^!}(-x)
 = 1 + \sum_{g \geq 1} \kappa^g \, Q_g(x),
\end{equation}
where $\kappa = \kappa(\cA)$ is the modular characteristic and
$Q_g(x)$ is a polynomial of degree $\leq 3g - 3$ encoding
genus-$g$ curvature corrections. At genus~$0$, this recovers
the classical relation $P_\cA(x) \cdot P_{\cA^!}(-x) = 1$.
\end{conjecture}

\noindent
This is the chiral replacement for the classical Koszul functional
equation $h_A(t) \cdot h_{A^!}(-t) = 1$, which fails chirally
because the OS algebra twists the bar complex by integration kernels
on $\FM_n(X)$. The genus correction $Q_g$ is a polynomial (not a
power series) because $\dim \mathcal{M}_{g,n}^{\mathrm{virt}} = 3g-3$
bounds the relevant strata.

\subsection{Fusion finiteness and bar rationality}

\begin{conjecture}[Fusion finiteness; \ClaimStatusConjectured]
\label{conj:fusion-finiteness-koszulness}
\index{fusion!finiteness}
\index{bar generating function!holonomic}
If $\cA$ is chirally Koszul with finitely many strong generators
of rank~$r$, then $P_\cA(x) = \sum_{n \geq 0}
\dim H^n(\barB(\cA))\, x^n$ is D-finite
\textup{(}holonomic: bar cohomology dimensions satisfy a
P-recursive recurrence of order $\leq r + 1$\textup{)}.
\end{conjecture}

\noindent
Koszulness implies $E_2$ collapse of the bar spectral sequence, which
forces algebraicity of degree $\leq 2$
(Conjecture~\ref{conj:bar-gf-algebraicity}). Algebraic functions
are D-finite but \emph{not} rational in general: the standard
examples $P_{\widehat{\mathfrak{sl}}_2}(x)$,
$P_{\mathrm{Vir}}(x)$, $P_{\beta\gamma}(x)$ all involve
$\sqrt{1 - 2x - 3x^2}$, which is algebraic of degree~$2$, not
rational. The correct characterization is \emph{diagonal purity}
($H^{n,k}(\barB(\cA)) = 0$ for $k \neq n + h$),
not mere rationality.
\label{conj:rationality-vs-koszulness}

\subsection{C$_2$-cofiniteness and the Koszul--finiteness interaction}
\label{conj:c2-cofiniteness-koszul-interaction}
\index{C2-cofiniteness@$C_2$-cofiniteness!interaction with Koszulness}

$C_2$-cofiniteness and Koszulness are \emph{orthogonal} conditions:
$V_k(\mathfrak{g})$ is always Koszul (universal PBW) but never
$C_2$-cofinite at generic~$k$; the simple quotient
$L_k(\mathfrak{g})$ at admissible level is $C_2$-cofinite but
Koszulness may fail (singular vectors source higher bar
differentials). Koszul duality exchanges the two: at
admissible~$k$ the associated variety
$X_{L_k} = \{0\}$, but at the dual level $k' = -k - 2h^\vee$ one
has $X_{L_{k'}} \neq \{0\}$ in general.

\begin{conjecture}[Koszul--$C_2$ duality; \ClaimStatusConjectured]
\label{conj:koszul-c2-duality}
\index{C2-cofiniteness@$C_2$-cofiniteness!Koszul duality}
\index{Koszul pair!C2 duality@$C_2$ duality}
At non-degenerate admissible level, $L_k(\mathfrak{g})$ satisfies:
\begin{enumerate}[label=\textup{(\alph*)},leftmargin=2em]
\item each fixed conformal-weight sector of $\barB(L_k)$ has
 finite-dimensional cohomology
 \textup{(}proved weightwise finiteness from
 $C_2$-cofiniteness\textup{)};
\item the character-theoretic bar spectral sequence is weightwise
 finite; on the worked $\widehat{\mathfrak{sl}}_2$ surface at
 $k=-1/2$ it yields the vacuum Kac--Wakimoto character formula,
 while on the general non-degenerate admissible lane it is compared
 against the external Kac--Wakimoto character formulas;
\item a stronger global statement \textup{(}finite-dimensional total
 bar cohomology, or a page bound by $\dim A(L_k)$\textup{)} is not
 part of the current proved surface.
\end{enumerate}
For associated varieties in type~$A$, Koszul duality implements
Barbasch--Vogan duality on nilpotent orbits
\textup{(}the involution $d \colon \mathcal{N}(\mathfrak{g})
\to \mathcal{N}(\mathfrak{g})$ sending a partition to its
transpose; see Barbasch--Vogan~\cite{BV85}\textup{)}:
\begin{equation}\label{eq:koszul-bv-duality}
 X_{L_{k'}(\mathfrak{g})} = \overline{d\bigl(X_{L_k(\mathfrak{g})}\bigr)},
 \qquad k' = -k - 2h^\vee.
\end{equation}
This is the conjectural type-$A$ simple-quotient pattern
\textup{(}Conjecture~\textup{\ref{conj:orbit-duality}}\textup{)};
the dual-level associated variety identification remains open.
Beyond type~$A$, the Barbasch--Vogan involution involves the
Spaltenstein map rather than partition transpose; the extension
to all types is open.
\end{conjecture}

\subsection{Moduli of Koszul parameters}

\begin{remark}[Moduli of Koszul parameters]
\label{rem:moduli-koszul-parameters}
\index{Koszul pair!moduli of parameters}
\index{admissible level!Koszul locus}
\index{critical level!fixed point}
The one-parameter family $V_k(\mathfrak{g})$ stratifies the
level line into three Koszul regimes.

\textup{(i)}
\emph{Generic level.}
The universal affine algebra $V_k(\mathfrak{g})$ has a PBW
filtration at every level, and its associated graded is the
polynomial chiral algebra on $\mathfrak g[t^{-1}]t^{-1}$.  On the
standard universal-affine surface this PBW input is the independent
Koszulness witness; it is not a consequence of the two-sided
Hochschild bar resolution
(Theorem~\ref{thm:chiral-bar-resolution-exact}).  At generic
non-admissible~$k$ one has $L_k=V_k$, so the simple quotient inherits
that universal-affine PBW/Koszul package.

\textup{(ii)}
\emph{Critical level} $k = -h^\vee$.
This is the unique fixed point of the Feigin--Frenkel
involution $k \mapsto -k - 2h^\vee$ on the \emph{level parameter}.
This fixed-point statement is not a theorem that the affine algebra
itself satisfies $\cA^! \cong \cA$ at critical level.
What is visible internally is narrower: the scalar curvature vanishes,
the universal affine algebra acquires a large center, and critical
bar-cobar inversion / fixed-point packages can be formulated on the
specific surfaces where they are separately proved. The screening-model
differential
$\mathrm{Sym}[S_1, \ldots, S_r] \otimes \Omega^*_{\log}$
is only a free-field model
\textup{(}Remark~\ref{rem:w-critical-bar-viewpoint}\textup{)},
not an identification used here.
The Feigin--Frenkel center
$\mathfrak{z}(\widehat{\mathfrak{g}}) \cong
\mathrm{Fun}(\mathrm{Op}_{{}^L\!\mathfrak{g}})$
supplies central operators and a Miura/screening model, while the
external critical-level Langlands viewpoint for $\mathcal W$-algebras
\textup{(}Remark~\ref{rem:frenkel-gaitsgory-langlands-viewpoint}\textup{)}
marks the comparison surface with Langlands-dual data.
This is distinct from the proved internal critical fixed-point
inversion package, which does not derive a Langlands-dual
identification from bar-cobar inversion alone.

\textup{(iii)}
\emph{Admissible levels.}
$L_k(\mathfrak{g})$ at $k \in \Sigma(\mathfrak{g})$
(Definition~\ref{def:admissible-level}) acquires singular
vectors that source higher bar differentials. Bar--cobar
inversion may fail; the Koszulness question reduces to whether
screening operators generate \emph{all} bar differentials
(Proposition~\ref{prop:bar-bgg-sl2}).
\end{remark}

\section{Open mathematical conditions}

The remaining problems are the following precise mathematical conditions.
\begin{enumerate}
\item Classify the chiral algebras for which the conilpotent bar
coalgebra $\Bbarch(\cA)$ resolves a dual coalgebra~$\cA^i$ and for
which the finite-type Verdier hypotheses produce
the algebra~$\cA^!$.
\item Construct the $\infty$-categorical module comparison attached
to a Koszul pair without replacing it by the bar-cobar counit:
the target is the factorization comodule category over
$\Bbarch(\cA)$, and only on the finite-type Verdier surface may it be
rewritten as a completed $\cA^!$-module category.
\item Extend the periodicity phenomena
(Theorems~\ref{thm:virasoro-periodicity}
and~\ref{thm:affine-periodicity-critical}) to non-simply-laced types
by computing the first unknown bar cohomology groups, beginning with
$H^2(\barB(\widehat{G}_2))$.
\end{enumerate}

\medskip

The Koszul pair structure developed in this chapter consists of the
quadratic condition, the PBW/completion replacements for
nonquadratic algebras, and the level-shifting involutions visible on
the standard families.  The examples of
Part~\ref{part:characteristic-datum} use this structure only after
the Koszul condition and the Verdier-dual lane have been identified
separately.  The chiral Hochschild complex
$Z^{\mathrm{der}}_{\mathrm{ch}}(\cA)$ controls deformations of the
bar complex in the proved Hochschild range; same-type periodicity is
transported across the finite-type module equivalence of
Corollary~\ref{cor:hochschild-ring-koszul}, while mixed-type
periodicity remains open.

\section{Synthesis}
\label{sec:kps-supplement}

\begin{remark}[K3 bar coalgebra and the Verdier-dual lane]
\label{rem:kps-1}
The K3 bialgebra $\mathbf H_{\Delta_5}$ supplies a Hall-theoretic
bar coalgebra
\[
C_{\mathrm{K3}} = B_{\mathrm{ch}}(\mathbf H_{\Delta_5})
\quad\text{on}\quad
(\mathrm{Ran}(E^{\mathrm{nod,sm}}_{24}),\; \bar{\mathcal A}_2).
\]
The cobar functor is not the second member of a Koszul pair; it is
the inversion functor applied to this coalgebra.  At the Lusztig
specialisation $\hbar^2=-1/8$ the Koszul resolution
\[
\mathrm{Bar}_\bullet(\mathbf H_{\Delta_5}\otimes \mathbf H_{\Delta_5}^{\mathrm{op}}, \mathbf H_{\Delta_5}, \mathbf H_{\Delta_5}) \;\xrightarrow{\;\simeq\;}\; \mathbf H_{\Delta_5}
\]
is a quasi-isomorphism and lifts to a strict isomorphism after
passing to the derived Hall category of $K3\times E$.  Thus the
bar-cobar counit gives
\[
 \Omega_X(C_{\mathrm{K3}})\xrightarrow{\sim}\mathbf H_{\Delta_5}.
\]
The Verdier-dual/Koszul algebra associated to this K3 surface is a
separate Vol.~III identification: one first Verdier-dualizes the Hall
bar coalgebra on the Ran--Siegel base and then applies the completed
cobar comparison under the stated convergence hypotheses.  The
statement is therefore an instance of bar-cobar inversion at the K3
base, not a definition of $\mathbf H_{\Delta_5}^!$ as
$\Omega_X B_{\mathrm{ch}}(\mathbf H_{\Delta_5})$. See Vol.~III,
Chapter~\ref{V3-ch:k3-chiral-bialgebra-platonic}, Theorem~A.
\end{remark}
